% Compile: pdflatex paper.tex && bibtex paper && pdflatex paper.tex && pdflatex paper.tex
\documentclass[11pt, letterpaper]{article}

\usepackage[margin=1in]{geometry}
\usepackage{amsmath,amssymb,amsthm}
\usepackage{booktabs}
\usepackage{xcolor}
\usepackage{microtype}
\usepackage{float}
\usepackage{caption}
\usepackage{graphicx}
\graphicspath{{../docs/out/}{figures/}}
\usepackage{tikz}
\usepackage[colorlinks=true,linkcolor=black,citecolor=black,urlcolor=blue!70!black,hypertexnames=false]{hyperref}
\usepackage{enumitem}
\usepackage{accsupp}
\usepackage{eso-pic}

\newcommand{\nocopytext}[1]{\BeginAccSupp{ActualText={}}#1\EndAccSupp{}}

% ── Preprint watermark ───────────────────────────────────────────────────
\AddToShipoutPictureBG{%
  \put(\LenToUnit{0.5\paperwidth},\LenToUnit{0.955\paperheight}){%
    \makebox(0,0){\nocopytext{\textcolor{gray!25}{\large\textit{Preprint --- not peer reviewed}}}}%
  }%
  \put(\LenToUnit{0.5\paperwidth},\LenToUnit{0.04\paperheight}){%
    \makebox(0,0){\nocopytext{\textcolor{gray!25}{\large\textit{Preprint --- not peer reviewed}}}}%
  }%
}

\usepackage{fancyhdr}
\fancypagestyle{plain}{%
  \fancyhf{}%
  \fancyfoot[C]{\nocopytext{\thepage}}%
  \renewcommand{\headrulewidth}{0pt}%
}
\pagestyle{plain}

% ── Theorem environments ────────────────────────────────────────────────
\newtheorem{theorem}{Theorem}[section]
\newtheorem{lemma}[theorem]{Lemma}
\newtheorem{proposition}[theorem]{Proposition}
\newtheorem{corollary}[theorem]{Corollary}
\newtheorem{conjecture}[theorem]{Conjecture}
\newtheorem{compresult}[theorem]{Computational Result}
\theoremstyle{definition}
\newtheorem{definition}[theorem]{Definition}
\newtheorem{example}[theorem]{Example}
\theoremstyle{remark}
\newtheorem{remark}[theorem]{Remark}

% ── Notation shortcuts ──────────────────────────────────────────────────
\newcommand{\Hom}{\mathrm{hom}}
\DeclareMathOperator{\Sp}{Sp}
\newcommand{\Aut}{\mathrm{Aut}}
\newcommand{\Wt}[1]{w_{#1}}
\newcommand{\ip}[2]{\langle #1, #2 \rangle_{\mathbf{s}}}

\title{%
  \textbf{Path Minimizers and Leontovich Thresholds in Tree Homomorphisms}\\[0.5em]
}
\author{%
  Shane Jaroch \\
  \small \textit{Department of Mathematics and Statistics,} \\
  \small \textit{School of Engineering and Computer Science} \\
  \small Oakland University \\
  \small \href{mailto:sjjaroch@oakland.edu}{\texttt{sjjaroch@oakland.edu}}
}
\date{\today}

\begin{document}
\maketitle

\begin{abstract}
We study graph homomorphisms from trees to fixed target graphs $H$, focusing
on Leontovich graphs: targets for which the path does not minimize the number
of homomorphisms. We prove that no $T(x,y,z)$ has a near-path crossover at
$n=5$. At $n=6$, for targets in this family, four of the six source-tree
comparisons are proved; the remaining two reduce to explicit polynomial
inequalities and remain open beyond the tested parameter range. The smallest
crossover threshold witnessed here is $n=7$. For symmetric trees with two or
three orbits, an odd near-path crossover can be decided by a single finite
sign check. For the four-orbit tree $T(7,1,9)$, we prove that its odd threshold
$n=13$ is permanent. We also prove that the bipartite double cover preserves
the Leontovich property. Our searches produce a 15-vertex witness $H^*$ by
exhaustive bipartite enumeration and an 18-vertex witness $H_{18}$ by SMT
synthesis. The latter first crosses at $(n,d)=(15,4)$ and is also minimal in
the completed depth-2 sweep through $n \le 51$. Finally, a
3,644-vertex double-cover witness first crosses at the even threshold
$n=17{,}340$, disproving the proposed bipartite parity hypothesis.
\end{abstract}

% ════════════════════════════════════════════════════════════════════════
\section{Introduction}
\label{sec:intro}

\subsection{Counting maps between graphs}

Let $G$ and $H$ be finite graphs, where $H$ may carry loops. A
\emph{homomorphism} $\varphi \colon G \to H$ is a map
$V(G) \to V(H)$ that sends edges to edges: if $uv \in E(G)$ then
$\varphi(u)\varphi(v) \in E(H)$. We write $\Hom(G,H)$ for the number of
such maps.

The quantity $\Hom(G,H)$ specialises to many familiar counts. If
$H = K_q$ is complete, it counts proper $q$-colorings; if $H$ is a
looped edge, it counts independent sets; if $H$ is a path, it counts
height functions with bounded increments. We call $H$ the
\emph{target}, and we often call $\Hom(G,H)$ the number of
\emph{$H$-colorings} of $G$.

Two extremal questions arise: among all trees $T$ on $n$ vertices,
which maximizes $\Hom(T,H)$, and which minimizes it? Sidorenko~\cite{sidorenko1991}
settled the first: the star $K_{1,n-1}$ is always a maximizer for every
target $H$. The second remains open in general.

\subsection{The path conjecture and its failure}

Galvin and Tetali~\cite{galvin-tetali} proposed the natural companion:
that the path $P_n$ should always be a minimizer. The intuition is
appealing. A homomorphism from a tree is built by propagating
constraints outward from a root, and the path is the tree that spreads
those constraints along the longest possible chain, giving each vertex
the fewest independent choices. Csikv\'ari and Lin~\cite{csikvari-lin}
developed the systematic tool for this circle of ideas, the
KC-transformation, and proved the conjecture for large classes of
targets.

The conjecture is false. Leontovich~\cite{leontovich1989} exhibited a
target for which some tree beats the path, and
Galvin, McMillon, Redlich and Nir~\cite{galvin2025} gave explicit
bipartite counterexamples. We call $H$ a \emph{Leontovich graph} when
the path fails to minimize, and we call the least such $n$ the
\emph{crossover threshold}. A companion paper~\cite{galvin2026}
introduces the dual class of \emph{Hoffman--London graphs}; here we ask
how small a Leontovich graph can be, and how early a crossover can
occur.

What makes the problem hard is that the failure is faint. In every known
example the counts differ by a relative margin of about $10^{-4}$ or
less, and the winner is usually just a single pendant move away from the
path. The signal is a near-cancellation of exponentially growing terms,
so crude estimates miss it and floating-point noise can erase it. Much
of what follows is therefore about comparisons that survive at that
scale, and about marking proofs separately from bounded searches.

\subsection{What this paper does}

We attack the problem in three ways.

\emph{Analytically.} We isolate an eigenvalue condition that reduces each
odd near-path comparison to a single base-case sign check
(Theorem~\ref{thm:spectral}), determine the exact odd crossover
threshold of the smallest previously known Leontovich tree
(Theorem~\ref{thm:threshold}), and show that the bipartite double cover
preserves both the Leontovich property and its asymptotic strengthening
(Theorem~\ref{thm:double_cover}).

\emph{Computationally.} We sweep large families of targets --- all
connected simple graphs on at most $11$ vertices, all connected
bipartite graphs on at most $15$ vertices, and all symmetric trees
below $78$ vertices --- against the near-path source trees defined
below, and we record the results as bounded sweeps rather than theorems.

\emph{By synthesis.} Where enumeration runs out, we encode the
crossover condition as an integer constraint problem and hand it to an
SMT solver. This produces the $18$-vertex depth-$2$ witness; the
$15$-vertex depth-dependent witness comes from exhaustive bipartite
enumeration.

\subsection{Standards of evidence}
\label{sec:evidence}

Because the arguments mix proof with computation, and because earlier
drafts contained claims that failed exact arithmetic, we label every
substantive assertion with one of four tags and use them consistently.

\begin{description}[leftmargin=2.4em,style=nextline,font=\normalfont\bfseries]
  \item[(P) Proved.] An unconditional proof, quantifiers included.
    Stated as \textbf{Theorem}, \textbf{Proposition}, \textbf{Lemma} or
    \textbf{Corollary}.
  \item[(E) Exact finite computation.] A numerical claim about named
    objects, verified in exact integer or certified
    arbitrary-precision arithmetic.
  \item[(S) Bounded sweep.] A negative result of the form ``no member
    of family $\mathcal{F}$ exhibits behaviour $B$ within the tested
    window''. This certifies the window, not the family; we never state it as a
    nonexistence theorem.
  \item[(N) High-precision numerics.] A floating-point computation with
    a residual-based error estimate but no outward-rounded interval or
    exact enclosure. This is numerical evidence, not a proof.
\end{description}

The distinction between (P) and (S) matters most. A near-path sweep
certifies only that the tested trees fail to beat the path; it does not
certify that no tree does. Accordingly, we reserve nonexistence claims
for proved results and state the scope of each sweep explicitly.

\subsection{Summary of contributions}

\begin{enumerate}[label=(\roman*)]
  \item \textbf{(E)} Verification of the path-minimizer theorem for the
    odd paths $H \in \{P_3,P_5,P_7,P_9\}$ over all source trees on
    $n \le 22$ vertices. For $P_3$ and $P_5$ the path is minimal but
    not unique; from $P_7$ onward it is uniquely minimal. The
    arithmetic mechanism is that, with $k$ the path-length parameter
    and $\lambda_1$ the first eigenvalue of the active quotient,
    $\lambda_1^2$ is rational exactly for
    $k \in \{1,2,3,5\}$. This resolves the question computationally
    through $n \le 22$ and identifies the arithmetic mechanism behind
    the observed transition in Problem~6.3 of~\cite{galvin2026}
    (Section~\ref{sec:odd-path}).
  \item \textbf{(P)} An eigenvalue criterion (Theorem~\ref{thm:spectral}):
    for a rooted bipartite target whose active quotient has a single
    positive eigenvalue, each fixed-depth odd near-path comparison has
    constant sign. For all $2$-orbit and $3$-orbit symmetric trees this
    reduces detection to one base-case check
    (Corollary~\ref{cor:2_and_3_orbit_trees}).
  \item \textbf{(P)} Among odd $n$, the crossover threshold of
    $T(7,1,9)$ is exactly $n = 13$, and it is permanent: the sign changes
    once and never again (Theorem~\ref{thm:threshold}). This threshold was already
    known in correspondence. The companion impossibility at $n = 5$ for
    the spherically symmetric family $T(x,y,z)$ is a polynomial
    identity (Proposition~\ref{prop:n5}).
  \item \textbf{(S)} All $1{,}018{,}680{,}329$ connected simple graphs on
    $m \le 11$ vertices, and all $608{,}930{,}559$ connected bipartite
    graphs on $m \le 15$ vertices, tested against the near-path family
    with $n \le 200$, $d \le 20$. No simple target below $12$ vertices
    produced a near-path crossover. This bounds, but does not settle,
    Problem~4.3 of~\cite{galvin2025} (Section~\ref{sec:bounds}).
  \item \textbf{(E)} Two small bipartite witnesses: a $15$-vertex graph
    $H^*$ whose earliest crossover is at $(n,d) = (49,16)$, and an
    $18$-vertex graph $H_{18}$ whose first near-path crossover is at
    $(n,d)=(15,4)$ with margin $\Delta=500{,}576$. It crosses against
    the classical depth-$2$ near-path at $n=17$ with margin
    $\Delta = 5{,}068{,}778$ (Section~\ref{sec:discussion}).
  \item \textbf{(P)} The bipartite double cover preserves both the
    Leontovich and the strongly Leontovich property
    (Theorem~\ref{thm:double_cover}). Applied to the looped tree
    $\hat T(1,35,1,50)$, this yields a \emph{simple} strongly Leontovich
    bipartite graph on $3{,}644$ vertices, whose first even crossover we
    locate exactly at $n = 17{,}340$ --- evidence that even-$n$
    crossovers do not require a loop, only enough room to delay the
    cancellation.
  \item \textbf{(N)} A perturbed double cover on $6{,}806$ vertices whose
    estimated leading-coefficient ratio is $0.999952140676 < 1$,
    a candidate simple non-bipartite strongly Leontovich target.
\end{enumerate}

Section~\ref{sec:prelim} fixes notation and records the two
transfer-matrix identities that every computation in the paper rests on.

\section{Preliminaries}
\label{sec:prelim}

Throughout, $H$ denotes a finite connected target graph, possibly with
loops, on $m = |V(H)|$ vertices, with adjacency matrix $A = A(H)$; a
loop at $u$ contributes $A_{uu} = 1$. Source graphs are trees on $n$
vertices. All logarithms are natural and all matrices are real.

\subsection{Walks are the only thing being counted}

Let $\mathbf{1}$ be the all-ones vector and set

\[
  w_k = A^k \mathbf{1},
  \qquad\text{so that}\qquad
  w_k(u) = \bigl|\{\text{walks of length } k \text{ in } H \text{ starting at } u\}\bigr|.
\]
A homomorphism from a tree is determined by choosing an image for a root
and then extending independently along each branch, so every count we
need is a sum of products of entries of the $w_k$.

\begin{lemma}[Path and spider counts]
  \label{lem:transfer}
  Let $\Sp(a_1,\dots,a_r)$ denote the \emph{spider}: the tree obtained by
  identifying one endpoint of $r$ paths with $a_1, \dots, a_r$ edges
  respectively at a common centre. A spider is a tree with one branch
  vertex; we write $\Sp(a,b,c)$ for the spider with leg lengths $a,b,c$.
  Root the spider at its centre. For each vertex $u \in V(H)$, count the
  homomorphisms that send the root to $u$; the legs then extend
  independently, so we sum over all such $u$.
  Then
  \[
    \Hom(P_n, H) \;=\; \mathbf{1}^{T} A^{\,n-1} \mathbf{1}
      \;=\; \sum_{u \in V(H)} w_{n-1}(u),
    \qquad
    \Hom\bigl(\Sp(a_1,\dots,a_r), H\bigr)
      \;=\; \sum_{u \in V(H)} \prod_{i=1}^{r} w_{a_i}(u).
  \]
\end{lemma}

\begin{proof}
  Root the spider at its centre. For each vertex $u \in V(H)$, count the
  homomorphisms that send the root to $u$; the branches are vertex-disjoint
  apart from the centre, so their images may be chosen independently. The
  choices along a branch with $a_i$ edges are exactly the walks of length
  $a_i$ from $u$. The path is the case $r = 2$, or directly the $k$-step
  iteration of $A$.
\end{proof}

\begin{figure}[H]
  \centering
  \begin{tikzpicture}[line cap=round, line join=round, every node/.style={font=\small}]
    \coordinate (root) at (0,0);
    \coordinate (leftend) at (-5.2,0);
    \coordinate (rightend) at (5.2,0);
    \coordinate (downend) at (0,-2.8);
    \draw[line width=0.8pt] (root) -- (leftend);
    \draw[line width=0.8pt] (root) -- (rightend);
    \draw[line width=0.8pt] (root) -- (downend);
    \foreach \x in {-4.2,-3.2,-2.2,-1.2,1.2,2.2,3.2,4.2} {
      \fill (\x,0) circle (1.8pt);
    }
    \foreach \y in {-0.9,-1.8,-2.7} {
      \fill (0,\y) circle (1.8pt);
    }
    \fill (root) circle (2.8pt);
    \node[above] at (0,0.45) {root};
    \node[left] at (-5.45,0) {$a$};
    \node[right] at (5.45,0) {$b$};
    \node[below] at (0,-3.15) {$c$};
    \node[below] at (0,-3.85) {spider $\,\Sp(a,b,c)$};
  \end{tikzpicture}
  \caption{A spider rooted at its centre.}
  \label{fig:spider-lemma}
\end{figure}

\subsection{The near-path family}

\begin{definition}[Near-path trees]
  \label{def:nearpath}
  For $n \ge 5$ and $1 \le d \le n-3$, the \emph{near-path tree} is the
  spider
  \[
    E_n^{(d)} \;=\; \Sp\bigl(n-d-2,\; d,\; 1\bigr),
  \]
  the tree on $n$ vertices obtained from the path $P_{n-1}$ by attaching
  a pendant vertex to the vertex at distance $d$ from one end. We write
  $E_n = E_n^{(2)}$ for the classical case, the tree used by
  Leontovich~\cite{leontovich1989}, obtained by hanging a pendant on the
  third-from-last vertex of $P_{n-1}$.
\end{definition}

Combining Definition~\ref{def:nearpath} with Lemma~\ref{lem:transfer},
\begin{equation}
  \label{eq:nearpath-count}
  \Hom\bigl(E_n^{(d)}, H\bigr)
    \;=\; \sum_{u \in V(H)} w_{n-d-2}(u)\, w_{d}(u)\, w_{1}(u),
\end{equation}
which is the identity implemented by every verification script in this
work. Two consequences are worth recording.

\begin{remark}[The depth parameter is a leg length]
  \label{rem:d-symmetry}
  A spider does not remember the order of its legs, so
  $E_n^{(d)} \cong E_n^{(n-2-d)}$. The parameter $d$ is therefore not a
  depth so much as the length of the shorter of the two long legs, and a
  sweep over $d \le d_{\max}$ covers exactly those near-path trees with
  $\min\{d,\, n-2-d\} \le d_{\max}$. Our sweeps use
  $d_{\max} = 20$ with $n \le 200$; the untested trees are those whose
  pendant sits near the middle of the path.
\end{remark}

\begin{remark}[Why this family and not another]
  \label{rem:why-nearpath}
  The near-path trees are the natural first place to look, for a reason
  made precise in Remark~\ref{rem:filter}: any tree that beats the path must
  match its exponential growth rate $\lambda_1^{\,n}$ and win on the
  constant in front. Branching in the interior raises the growth rate
  outright, so a challenger must confine its defect to the boundary.
  This is a heuristic for choosing a search family, not a theorem, and
  we treat every negative result obtained with it accordingly.
\end{remark}

\subsection{Symmetric trees and their quotients}

\begin{definition}[Spherically symmetric trees]
  \label{def:symtree}
  For $d_1,\dots,d_k \ge 1$, the tree $T(d_1,\dots,d_k)$ is rooted, of
  depth $k$, in which every vertex at depth $i-1$ has exactly $d_i$
  children. Its vertex set falls into $k+1$ \emph{orbits} by depth, of
  sizes $s_0 = 1$ and $s_i = d_1 d_2 \cdots d_i$, and
  \[
    |V| \;=\; 1 + d_1 + d_1 d_2 + \cdots + d_1 d_2 \cdots d_k .
  \]
  We write $\hat T(d_1,\dots,d_k)$ for the same tree with a loop added at
  the root. The case $k = 3$ recovers the classical family
  $T(x,y,z)$ on $1 + x + xy + xyz$ vertices, whose four orbits are root,
  children, grandchildren and leaves.
\end{definition}

The loop matters: it destroys bipartiteness, and with it the parity
obstruction of Section~\ref{sec:leontovich}.

\begin{definition}[Equitable partition and quotient matrix]
  \label{def:quotient}
  A partition $\pi = \{C_1,\dots,C_K\}$ of $V(H)$ is \emph{equitable} if
  there are integers $q_{ij}$ such that every vertex of $C_i$ has exactly
  $q_{ij}$ neighbours in $C_j$. The matrix $Q = (q_{ij})$ is the
  \emph{quotient matrix} of $\pi$ and $\mathbf{s} = (|C_1|,\dots,|C_K|)$
  its \emph{size vector}. The orbit partition of any automorphism group
  is equitable; so is the coarsest partition produced by color refinement.
\end{definition}

\begin{lemma}[Reduction to the quotient]
  \label{lem:quotient}
  Let $\pi$ be equitable with quotient $Q$ and size vector $\mathbf{s}$.
  Then each walk vector $w_k$ is constant on cells, and its cell values
  are the entries of $Q^k \mathbf{1}_K$. Consequently, for every $n$ and
  every spider,
  \[
    \Hom(P_n,H) \;=\; \mathbf{s}^{T} Q^{\,n-1}\mathbf{1}_K,
    \qquad
    \Hom\bigl(\Sp(a_1,\dots,a_r),H\bigr)
      \;=\; \sum_{i=1}^{K} s_i \prod_{j=1}^{r} \bigl(Q^{a_j}\mathbf{1}_K\bigr)_i .
  \]
\end{lemma}

\begin{proof}
  Induct on $k$. The claim holds for $w_0 = \mathbf{1}$. If $w_{k}$ takes
  the value $x_i$ on $C_i$, then for $u \in C_i$,
  $w_{k+1}(u) = \sum_j q_{ij} x_j$, which depends only on $i$. Summing
  Lemma~\ref{lem:transfer} cell by cell gives the two displayed
  identities.
\end{proof}

Lemma~\ref{lem:quotient} is what makes the computations feasible: a
target on $6{,}806$ vertices becomes an $18 \times 18$ integer matrix
with no loss of information about $\Hom(T,\cdot)$ for the trees we test.
Only the eigenvalues of $Q$ can appear in these counts; the rest of
$A(H)$ is invisible here, and we call those eigenvalues \emph{inert}.

\begin{lemma}[Two-step recurrence from a symmetric active spectrum]
  \label{lem:active_recurrence}
  If the active spectrum of $Q$ is contained in
  $\{\lambda_1, 0, -\lambda_1\}$, then the quotient-walk vectors satisfy
  \[
    Q^{k+2}\mathbf{1}_K \;=\; \lambda_1^{\,2}\, Q^k\mathbf{1}_K
  \]
  for every $k \ge 1$. Equivalently, after the kernel component has
  vanished, the walk counts obey $w_{k+2} = \lambda_1^2 w_k$.
\end{lemma}

\begin{proof}
  Expand $\mathbf{1}_K$ in the $\langle\cdot,\cdot\rangle_{\mathbf{s}}$-
  orthonormal eigenbasis of $Q$. The zero-eigenspace contributes nothing
  once $k \ge 1$, and the $\pm \lambda_1$ terms pick up the same factor
  $\lambda_1^2$ after two more steps.
\end{proof}

\begin{lemma}[The quotient is self-adjoint in the size weighting]
  \label{lem:selfadjoint}
  Define $\ip{x}{y} = \sum_i s_i x_i y_i$. Then $Q$ is self-adjoint with
  respect to $\langle\cdot,\cdot\rangle_{\mathbf{s}}$; in particular its
  eigenvalues are real and it has a $\langle\cdot,\cdot\rangle_{\mathbf{s}}$-orthonormal
  eigenbasis.
\end{lemma}

\begin{proof}
  Counting edges between $C_i$ and $C_j$ in two ways gives
  $s_i q_{ij} = s_j q_{ji}$, i.e.\ $\ip{Qx}{y} = \ip{x}{Qy}$.
  Equivalently $S^{1/2} Q S^{-1/2}$ is a symmetric matrix, where
  $S = \operatorname{diag}(\mathbf{s})$.
\end{proof}

\subsection{Leading coefficients and the two Leontovich properties}

Let $\lambda_1 \ge \lambda_2 \ge \cdots \ge \lambda_K$ be the eigenvalues
of $Q$ with $\langle\cdot,\cdot\rangle_{\mathbf{s}}$-orthonormal
eigenvectors $u^{(1)},\dots,u^{(K)}$, and put
$\alpha_j = \ip{\mathbf{1}}{u^{(j)}}$. Lemma~\ref{lem:quotient} and
Lemma~\ref{lem:selfadjoint} give the exact expansions
\begin{equation}
  \label{eq:expansion}
  \Hom(P_n,H) = \sum_{j=1}^{K} \alpha_j^{2}\, \lambda_j^{\,n-1},
  \qquad
  \Hom\bigl(E_n^{(d)},H\bigr)
    = \sum_{j=1}^{K} \alpha_j \beta_j^{(d)}\, \lambda_j^{\,n-d-2},
  \qquad
  \beta_j^{(d)} = \ip{w_1 \odot w_d}{u^{(j)}},
\end{equation}
where $\odot$ is the entrywise product. Everything in this paper is a
statement about the interference between the $j = 1$ term and the rest.

\begin{definition}[Leontovich, strongly Leontovich]
  \label{def:leontovich}
  A target $H$ is \emph{Leontovich} if there is an $n$ and a tree $T$ on
  $n$ vertices with $\Hom(T,H) < \Hom(P_n,H)$: the path is not a
  minimizer. It is \emph{strongly Leontovich}~\cite{galvin2026} if
  moreover a fixed near-path family beats the path for all large $n$,
  i.e.\ $\Hom(E_n^{(d)},H) < \Hom(P_n,H)$ for all sufficiently large $n$
  of the relevant parity.
\end{definition}

If $H$ is connected and non-bipartite, $\lambda_1$ strictly dominates
$|\lambda_j|$ for $j \ge 2$, and \eqref{eq:expansion} identifies the
strong property with a single inequality between leading coefficients.

\begin{definition}[Leading-coefficient ratio]
  \label{def:rho}
  For a connected non-bipartite target $H$ with active quotient
  $(Q,\mathbf{s})$, Perron eigenvalue $\lambda_1$ and positive Perron
  eigenvector $u = u^{(1)}$, set
  \begin{equation}
    \label{eq:rho}
    \rho_d(H) \;=\;
      \lim_{n \to \infty}
      \frac{\Hom\bigl(E_n^{(d)},H\bigr)}{\Hom(P_n,H)}
    \;=\;
      \frac{\ip{w_1 \odot w_d}{u}}
           {\lambda_1^{\,d+1}\, \ip{\mathbf{1}}{u}} .
  \end{equation}
  The right-hand side is independent of the normalisation of $u$. Then
  $H$ is strongly Leontovich for depth $d$ precisely when
  $\rho_d(H) < 1$, and $\rho_d(H) > 1$ forbids it.
\end{definition}

\begin{remark}[Bipartite targets and parity]
  \label{rem:bipartite-parity}
  If $H$ is bipartite then $-\lambda_1$ is also an eigenvalue of $Q$ and
  the limit in \eqref{eq:rho} need not exist: the two dominant modes
  reinforce at one parity of $n$ and cancel at the other. This is not a
  technicality but the governing phenomenon. It is why every crossover
  in this paper occurs at odd $n$ unless the target is very large
  (Section~\ref{sec:doublecover}), and why the looped root in
  $\hat T(\cdot)$ is the standard device for removing the obstruction.
  For bipartite $H$ we therefore read \eqref{eq:rho} along a fixed
  parity class and say so explicitly.
\end{remark}

\begin{remark}[The margins involved]
  \label{rem:margins}
  To fix expectations: the smallest witnesses in this paper win by a
  relative margin of about $4.5 \times 10^{-4}$ ($H_{18}$ at $(n,d)=(15,4)$)
  down to $4.8 \times 10^{-5}$ (the perturbed cover of
  Section~\ref{sec:doublecover}). A double-precision evaluation of
  \eqref{eq:nearpath-count} accumulates relative error on the order of
  $n\,m\,\varepsilon \approx 10^{-12}$, which is comfortably below these
  margins but not below every margin that might exist. Floating point is
  therefore used only to \emph{flag} candidates; every reported witness
  is re-derived in exact integer arithmetic through
  Lemma~\ref{lem:quotient}, and every asymptotic claim is accompanied by
  the high-precision error estimate described in
  Section~\ref{sec:evidence}.
\end{remark}

\subsection{Notation summary}

\begin{table}[H]
\centering
\caption{Notation used throughout.}
\label{tab:notation}
\begin{tabular}{@{}ll@{}}
  \toprule
  Symbol & Meaning \\
  \midrule
  $\Hom(G,H)$ & number of homomorphisms $G \to H$ \\
  $m$, $n$ & order of the target $H$; order of the source tree \\
  $A$, $w_k = A^k\mathbf{1}$ & adjacency matrix of $H$; walk-count vector \\
  $P_n$ & path on $n$ vertices ($n-1$ edges) \\
  $\Sp(a_1,\dots,a_r)$ & spider with legs of $a_1,\dots,a_r$ edges \\
  $E_n^{(d)} = \Sp(n{-}d{-}2, d, 1)$ & near-path tree; $E_n = E_n^{(2)}$ \\
  $T(d_1,\dots,d_k)$, $\hat T(\cdot)$ & symmetric tree; same with a looped root \\
  $Q$, $\mathbf{s}$, $K$ & quotient matrix, cell sizes, number of cells \\
  $\lambda_1 > |\lambda_j|$, $u$ & Perron eigenvalue and eigenvector of $Q$ \\
  $\lambda_2^{*}$ & largest active subdominant eigenvalue \\
  $\rho_d(H)$ & leading-coefficient ratio, \eqref{eq:rho} \\
  $\Delta_n = \Hom(P_n,H) - \Hom(E_n^{(d)},H)$ & crossover margin; $\Delta_n > 0$ means the path loses \\
  \bottomrule
\end{tabular}
\end{table}

\section{The Path Minimizer Theorem}
\label{sec:odd-path}

\begin{theorem}[Path Minimizer, Csikv\'ari and Lin \cite{csikvari2015}]
  For every odd path $P_k$ and every $n$-vertex tree $T$,
  $\Hom(P_n, P_k) \le \Hom(T, P_k)$.
\end{theorem}

\subsection{Computational verification}

To investigate the uniqueness of the path minimizer, we computationally
evaluated homomorphism counts for $H \in \{P_3, P_5, P_7, P_9\}$ and all
$n$-vertex trees with $n \le 22$. Tree enumeration uses the
Beyer--Hedetniemi canonical level-sequence algorithm~\cite{BeyerH80} implemented in C++,
producing all non-isomorphic trees. Tree counts are cross-validated
against OEIS A000055~\cite{oeis_A000055}.

\subsection{Regimes of minimality}

The eigenvalues of $P_k$ govern the degree of uniqueness with
which $P_n$ minimizes homomorphisms. We identify three regimes.

\subsubsection{Structural degeneracy}

\begin{proposition}
  \label{prop:p3}
  For any tree $T$ with bipartition $(A, B)$,
  \[
    \Hom(T, P_3) = 2^{|A|} + 2^{|B|}.
  \]
  In particular, $\Hom(P_n, P_3) = 2^{\lceil n/2 \rceil}
  + 2^{\lfloor n/2 \rfloor}$.
\end{proposition}

Since $P_n$ has the most balanced bipartition among $n$-vertex trees,
it achieves the minimum, but any tree sharing the same bipartition sizes
ties it exactly. Thus, the path minimizes the homomorphism count, but
it does not do so uniquely. This is a special case of Theorem~2.6
of~\cite{galvin2026}, since $P_3 \cong K_{1,2}$.

\subsubsection{\texorpdfstring{$P_5$}{P5}: Interior degeneracy and 2-shift invariance}

The adjacency matrix $A = A_{P_5}$ satisfies
\begin{equation}
  \label{eq:shift}
  A^3 \mathbf{1} = 3\, A \mathbf{1},
\end{equation}
where $\mathbf{1}$ is the all-ones vector. This identity implies that
the cube of the transfer matrix acts as a scalar multiple of the
matrix itself on the constant vector. Combinatorially, moving a pendant
edge two positions along the interior of a path preserves the
homomorphism count into $P_5$: the number of walks of length 3
originating from any vertex $v$ is exactly three times the number of
length-1 walks originating from $v$. Thus pendant relocations that
shift by 2 steps produce identical walk sums. This generates
exponentially growing tie classes ($\sim\!200$ ties at $n = 20$).

The eigenvalues of $P_5$ are $\pm\sqrt{3}, \pm 1, 0$, so
$\lambda_1^2 = 3$ exactly. In the language of
Lemma~\ref{lem:active_recurrence}, the active spectrum is
$\{\sqrt{3}, 0, -\sqrt{3}\}$ and the walk counts satisfy
$w_{k+2} = 3w_k$ for $k \ge 1$. The ratio
$\Hom(P_{n+2}, P_5) / \allowbreak \Hom(P_n, P_5) = 3$ for all odd $n$.

\subsubsection{Unique minimality}

For $k \ge 7$, the eigenvalues of $A_{P_k}$ are
$2\cos\bigl(\frac{j\pi}{k+1}\bigr)$ for $j = 1, \ldots, k$.
The squared spectral radius is $\lambda_1^2 = 4\cos^2\bigl(\frac{\pi}{k+1}\bigr) = 2 + 2\cos\bigl(\frac{2\pi}{k+1}\bigr)$.
By Niven's Theorem, $\cos(q\pi)$ is rational for rational $q$ if and only if
$q \in \{0, \pm 1/3, \pm 1/2, \pm 2/3, \pm 1\}$. Thus, $\lambda_1^2 \in \mathbb{Q}$ if and only if $k \in \{1, 2, 3, 5\}$.

For all $k \ge 6$, the squared dominant eigenvalue $\lambda_1^2$ is
irrational. This irrationality breaks any potential integer
shift-symmetry of the form $A^m \mathbf{1} = cA \mathbf{1}$ (which
would require $\lambda_1^m \in \mathbb{Q}$ or $\lambda_1^{m-1} \in
\mathbb{Q}$). In the computations above, this is exactly the mechanism
that separates the $P_5$ tie classes from the verified unique
minimality of $P_n$ for odd $n \ge 7$.

\begin{remark}[Irrationality]
  The verified transition to unique minimality occurs at $P_7$, the
  first odd path whose squared spectral radius ($2+\sqrt{2}$) is
  irrational. While still constructible in real radicals, that
  irrationality prevents the exact degeneracy classes seen in $P_5$.
\end{remark}

% ════════════════════════════════════════════════════════════════════════

\section{Analytical Properties of Leontovich Graphs}
\label{sec:leontovich}

\subsection{An eigenvalue condition that forbids crossovers}

\begin{theorem}[Single Positive Eigenvalue Reduction]
  \label{thm:spectral}
  Let $H$ be a rooted bipartite graph whose automorphism orbits respect
  the bipartition (i.e., no orbit contains vertices from both sides).
  If the quotient matrix of $H$ has exactly one positive eigenvalue,
  then the odd-$n$ near-path comparison has constant sign: for each fixed
  depth~$d$, the sign of $\Hom(P_n,H)-\Hom(E_n^{(d)},H)$ is determined on
  all odd~$n$ by any single verified odd base case.
\end{theorem}

\begin{proof}
  Since the orbits respect the bipartition, the quotient matrix has
  block form $\left[\begin{smallmatrix}0&B\\B^T&0\end{smallmatrix}\right]$
  after ordering the cells by side, so its nonzero spectrum is symmetric
  about zero.
  Because the base vector $\mathbf{1}$ and the initial degree vector
  are constant on automorphism orbits, the entire dynamic programming
  sequence $\{A^k \mathbf{1}\}$ remains in the column space of the
  quotient matrix, so the full adjacency spectrum is unnecessary.
  If there is exactly one positive eigenvalue $\lambda_1$, then every
  other nonzero eigenvalue is $-\lambda_1$, and the active spectrum is
  contained in $\{\lambda_1, 0, -\lambda_1\}$. By the quotient-eigenbasis
  expansion and Lemma~\ref{lem:active_recurrence},
  \[
    \Hom(G_n, H) = c\,\lambda_1^n + d\,(-\lambda_1)^n.
  \]
  For odd $n$, this reduces to $(c - d)\,\lambda_1^n$, and the kernel
  component vanishes once the exponent is at least $1$. There are no
  secondary exponential terms, so the ratio
  $\Hom(E_n, H) / \Hom(P_n, H)$ is constant on odd $n$. Thus the
  eigenvalue condition reduces any odd near-path crossover question to
  the sign of one verified odd base case in the family being studied.
\end{proof}

\begin{corollary}
  \label{cor:2_and_3_orbit_trees}
  All 2-orbit trees (stars $S_m$) and 3-orbit trees $T(x,y)$ have
  exactly one positive quotient eigenvalue. For these families, the
  eigenvalue condition reduces odd near-path crossover detection to a
  single base-case sign check.
\end{corollary}

\begin{corollary}[Double-cover transfer for small sweeps]
  \label{cor:loopy_transfer}
  If no connected simple graph on at most $2m$ vertices is Leontovich,
  then no connected graph on at most $m$ vertices, loops allowed, is
  Leontovich. In particular, if a connected looped target $H$ is
  non-bipartite, then its connected double cover $H \times K_2$ converts
  any unconditional simple-graph nonexistence bound at order~$\le 2m$ into
  the corresponding looped bound at order~$\le m$. Our bounded near-path
  sweeps therefore transfer only within their tested ranges: the
  simple-graph sweep up to $11$ vertices gives looped-target evidence on at
  most $5$ vertices for $n \le 200$, $d \le 20$, and the bipartite sweep
  up to $14$ vertices gives the analogous bounded evidence up to
  $7$ looped vertices.
\end{corollary}

\subsection{Four-orbit trees}

For 4-orbit trees $T(x,y,z)$, the characteristic polynomial of the
quotient matrix is
\begin{equation}
  \label{eq:charpoly}
  \lambda^4 - (x+y+z)\,\lambda^2 + xz = 0.
\end{equation}
Setting $s = x + y + z$, the discriminant of the quadratic in
$\lambda^2$ is
\[
  s^2 - 4xz = (x - z)^2 + y^2 + 2y(x + z).
\]
Since $x, y, z \ge 1$, this expression is always strictly positive
($\ge 5$). By Descartes' rule of signs, both roots for $\lambda^2$ are
positive, so \emph{every} spherically symmetric 4-orbit tree possesses
exactly 2 positive eigenvalues.

\subsection{The crossover threshold for \texorpdfstring{$T(7,1,9)$}{T(7,1,9)}}

This result was known privately to JD Nir and David Galvin prior to this work;
we record here the first fully written proof.

\begin{theorem}[Permanent Crossover]
  \label{thm:threshold}
  Let $H = T(7,1,9)$. Then $\Hom(E_n, H) < \Hom(P_n, H)$ for all
  odd $n \ge 13$, and direct computation gives $\Delta(11) < 0$ and
  $\Delta(13) > 0$.
\end{theorem}

\begin{proof}
  The characteristic polynomial of the quotient matrix
  \[
    M = \begin{pmatrix} 0&7&0&0 \\ 1&0&1&0 \\
    0&1&0&9 \\ 0&0&1&0 \end{pmatrix}
  \]
  is $t^4 - 17t^2 + 63 = 0$. Setting $u = t^2$:
  \[
    u = \frac{17 \pm \sqrt{37}}{2},
    \qquad
    \lambda_1 = \sqrt{\frac{17 + \sqrt{37}}{2}} \approx 3.3973,
    \qquad
    \lambda_2 = \sqrt{\frac{17 - \sqrt{37}}{2}} \approx 2.3364.
  \]
  For odd $n$, the quotient-eigenbasis expansion yields
  \[
    \Delta_{\mathrm{odd}}(n) \;=\; \Hom(P_n, H) - \Hom(E_n, H)
    \;=\; A\,\lambda_1^n - |B|\,\lambda_2^n,
  \]
  for constants $A > 0$ and $B \neq 0$. Setting
  $r = \lambda_2 / \lambda_1 = \sqrt{(17 - \sqrt{37})/(17 + \sqrt{37})}
  \approx 0.6877$:
  \[
    \Delta_{\mathrm{odd}}(n) = \lambda_1^n \bigl(A - |B|\,r^n\bigr).
  \]
  Since $r^n$ is strictly monotonically decreasing, $A - |B|\,r^n$ can
  cross zero at most once. Evaluating directly:
  \begin{align*}
    \Delta(11) &= 13{,}993{,}266 - 14{,}000{,}700 = -7{,}434 < 0, \\
    \Delta(13) &= 161{,}209{,}734 - 161{,}161{,}476 = +48{,}258 > 0.
  \end{align*}
  Thus $n = 13$ is the unique, permanent threshold.
\end{proof}

\begin{remark}[Parity heuristic, not theorem]
  \label{rem:parity}
  For even $n$, $P_n$ has a perfectly balanced bipartition $(n/2, n/2)$,
  minimizing its dominant projection. For even depths $d$ (e.g., $d=2$)
  where $E_n^{(d)}$ remains balanced, the path's linear geometry yields
  a smaller projection onto the dominant eigenvector. For odd depths
  (e.g., $d=1$), $E_n^{(d)}$ is forced into an unbalanced
  $(n/2+1, n/2-1)$ bipartition, giving it a strictly larger leading
  coefficient. This mechanism explains why even-$n$ crossovers were absent
  from many small sweeps, but Section~\ref{sec:doublecover} shows that it
  does not yield a universal obstruction. For odd $n$, both $P_n$ and $E_n^{(d)}$ share the same
  unbalanced bipartition $(\lceil n/2 \rceil, \lfloor n/2 \rfloor)$,
  enabling the secondary eigenvector to interfere and drive the crossover.
\end{remark}

\subsection{Minimum size in the symmetric family}

Exhaustive grid search over $T(x,y,z)$ with
$|V| = 1 + x + xy + xyz \le 100$ confirms that $T(7,1,9)$ with 78
vertices is the smallest Leontovich graph found in the $T(x,y,z)$
family within the tested range. A generalized search over all spherically symmetric trees
$T(d_1, \ldots, d_k)$ with $|V| \le 77$ and depth $k \le 6$ tested
5,732 trees with exact integer arithmetic
(\texttt{scripts/verify\_leontovich.py}, $n \le 200$, $d \le 20$) and
confirmed zero Leontovich instances (Table~\ref{tab:symmetric}).

\begin{table}[H]
\centering
\caption{Generalized symmetric tree search ($|V| \le 77$, exact integer
arithmetic, $n \le 200$, $d \le 20$).}
\label{tab:symmetric}
\begin{tabular}{@{}cccc@{}}
  \toprule
  Depth $k$ & Orbits & Trees tested & Leontovich found \\
  \midrule
  1 & 2 &    76 & \textbf{0} \\
  2 & 3 &   268 & \textbf{0} \\
  3 & 4 &   592 & \textbf{0} \\
  4 & 5 & 1,036 & \textbf{0} \\
  5 & 6 & 1,573 & \textbf{0} \\
  6 & 7 & 2,187 & \textbf{0} \\
  \midrule
  \multicolumn{2}{c}{Total} & 5,732 & \textbf{0} \\
  \bottomrule
\end{tabular}
\end{table}

\subsection{The \texorpdfstring{$n = 5$}{n=5} impossibility}

\begin{proposition}[$n = 5$ threshold is algebraically impossible]
  \label{prop:n5}
  For all $T(x,y,z)$ with $x, y, z \ge 1$,
  \[
    \Hom(P_5, T) - \Hom(E_5, T) = -xy\bigl((x - z - 1)^2 + y^2 z\bigr) < 0.
  \]
  In particular, $E_5$ always has \emph{more} homomorphisms than $P_5$
  into any $T(x,y,z)$, so no Leontovich crossover can occur at $n = 5$.
\end{proposition}

\begin{proof}
  Let $M$ denote the $4 \times 4$ quotient matrix of $T(x,y,z)$
  over its four automorphism orbits (root, children, grandchildren,
  leaves), with orbit-size vector
  $\mathbf{s} = (1,\, x,\, xy,\, xyz)^T$:
  \[
    M = \begin{pmatrix} 0&x&0&0 \\ 1&0&y&0 \\
    0&1&0&z \\ 0&0&1&0 \end{pmatrix}.
  \]
  For any path $P_n$, the homomorphism count decomposes as
  $\Hom(P_n, H) = \mathbf{s}^T M^{n-1} \mathbf{1}$,
  where $\mathbf{1} = (1,1,1,1)^T$. Iterating:
  \begin{align*}
    M\,\mathbf{1} &= (x,\; y+1,\; z+1,\; 1)^T, \\
    M^2\mathbf{1} &= (xy+x,\; x+yz+y,\; y+z+1,\; z+1)^T.
  \end{align*}
  The tree $E_5$ has vertices $\{1,2,3,4,5\}$ with edges
  $(1,2),\,(2,3),\,(3,4),\,(2,5)$, so vertex~2 has degree~3. Rooting
  at vertex~2 and summing over orbit assignments:
  \[
    \Hom(E_5, H) = \sum_{b=0}^{3}
      s_b \cdot (M\mathbf{1})_b^2 \cdot (M^2\mathbf{1})_b,
  \]
  since the three branches from vertex~2 contribute
  one factor of $(M\mathbf{1})_b$ each for the two pendant
  leaves (vertices~1 and~5) and one factor of
  $(M^2\mathbf{1})_b$ for the length-2 branch (vertices 3--4).

  Expanding $\mathbf{s}^T M^4 \mathbf{1}$ and collecting terms:
  \begin{align*}
    \Hom(P_5, T) - \Hom(E_5, T)
    &= -xy\bigl(x^2 - 2x(z+1) + (z+1)^2 + y^2 z\bigr) \\
    &= -xy\bigl((x - z - 1)^2 + y^2 z\bigr).
  \end{align*}
  Since $x, y, z \ge 1$, the squared term
  $(x - z - 1)^2 \ge 0$ and $y^2 z \ge 1$, so the
  parenthesised sum is strictly positive. Multiplied by $-xy < 0$,
  the difference is strictly negative.
\end{proof}

\begin{remark}[The remaining $n=6$ cases]
  \label{rem:n6_residue}
  The six trees on six vertices split into four proved cases and two
  checked cases. The path is the reference case; the star $K_{1,5}$
  follows from Sidorenko's theorem; the balanced near-path $\Sp(2,2,1)$
  is controlled by Theorem~\ref{thm:spectral}; and the double star
  $S_{2,2}$ is covered by the Cauchy--Schwarz bound above.
  The remaining two six-vertex trees, $\Sp(3,1,1)$ and
  $\Sp(2,1,1,1)$, are supported only by exact computation. Consequently,
  the four proved cases carry label (P), while these two cases carry label
  (E). A proof covering all six cases would require resolving this gap; the
  near-path sufficiency conjecture also remains open.
\end{remark}

\subsection{Minimum threshold landscape}

Sweeping the generator over all $T(x,y,z)$ with $|V| \le 500$ and
selecting the smallest graph achieving each odd-$n$ threshold yields
Table~\ref{tab:thresholds}. The vertex count bottoms out at $n = 13$
with $T(7,1,9)$, then \emph{increases} again---the threshold at $n = 15$
requires a larger graph than $n = 13$.

\begin{table}[H]
\centering
\caption{Smallest Leontovich graph by crossover threshold within the tested
range.}
\label{tab:thresholds}
\begin{tabular}{@{}llcc@{}}
  \toprule
  Family & Parameters & $|V|$ & Threshold \\
  \midrule
  $T(x,y,z)$ & $T(16, 1, 19)$ & 337 & $n = 7$ \\
              & $T(9, 1, 11)$  & 118 & $n = 9$ \\
              & $T(8, 1, 10)$  &  97 & $n = 11$ \\
              & $T(7, 1, 9)$   &  78 & $n = 13$ \\
              & $T(8, 1, 12)$  & 113 & $n = 15$ \\
  \bottomrule
\end{tabular}
\end{table}

\begin{remark}
  Galvin et al.~\cite{galvin2025} used $T(18,3,32)$ with $1{,}801$ vertices to prove
  the $n = 7$ crossover (their Theorem~1.3). Our landscape sweep shows
  that $T(16,1,19)$ achieves the identical crossover with only
  $337$~vertices---a $5.3\times$ reduction.  This is the smallest
  $T(x,y,z)$ graph achieving a threshold $n \le 7$ within the tested range.
\end{remark}

% ════════════════════════════════════════════════════════════════════════

\section{Bounding the Smallest Leontovich Graph}
\label{sec:bounds}

\subsection{Exhaustive sweep}

We tested all 1,018,690,329 connected simple graphs on $m \le 11$
vertices~\cite{oeis_A001349} (generated by McKay's
\texttt{geng}~\cite{mckay2014practical}) against the
$E_n^{(d)}$ filter for $n \le 200$, $d \le 20$.

\begin{compresult}[Computational sweep certificate]
  No connected simple graph on $m \le 11$ vertices produced a near-path
  violation $\Hom(E_n^{(d)},H)<\Hom(P_n,H)$ in the tested range
  $n \le 200$, $d \le 20$. For $m \le 9$, this was supplemented by
  brute-force checks over all source trees on $n \le 15$ vertices.
\end{compresult}

This computationally bounds the tested near-path obstruction in
Problem~4.3 of Galvin et al.~\cite{galvin2025}; it should not be read as an
unconditional proof that no other source-tree shape can beat the path for
$m=10,11$. Because our search space comprised simple graphs (generated by
\texttt{geng}), this evidence also does not exclude the possibility of a
looped or general graph on fewer than 12 vertices being strongly Leontovich
in the sense of~\cite{galvin2026}.
Galvin et al.~\cite{galvin2025} established an upper bound of
$|V| = 961{,}601$ using 4-orbit symmetric trees with a single looped
root. Expanding the degrees of freedom to 5-orbit looped symmetric trees
gives a smaller verified strongly Leontovich example on $1{,}822$ vertices
at $\hat{T}(1, 35, 1, 50)$, which will be used below because its stronger
margin is convenient for the double-cover construction. Furthermore, since our exhaustive search
includes all simple non-bipartite graphs (those containing odd cycles), this
sweep gives lower-bound evidence for simple strongly Leontovich graphs within
the tested near-path regime, significantly narrowing the uncertainty gap from
the bottom up. It cannot be upgraded directly to an unconditional small-target
nonexistence theorem: Section~\ref{sec:doublecover} exhibits a first crossover
only at $n = 17{,}340$, far beyond the sweep horizon used here.

\begin{table}[H]
\centering
\caption{Exhaustive sweep: all connected $H$ with $|V(H)| \le 11$.}
\label{tab:exhaustive}
\begin{tabular}{@{}ccccc@{}}
  \toprule
  $m$ & Connected $H$ & Cumulative & Operations & Violations \\
  \midrule
  1 & 1       & 1       & ---                     & \textbf{0} \\
  2 & 1       & 2       & ---                     & \textbf{0} \\
  3 & 2       & 4       & ---                     & \textbf{0} \\
  4 & 6       & 10      & $1.8 \times 10^7$   & \textbf{0} \\
  5 & 21      & 31      & $9.9 \times 10^7$   & \textbf{0} \\
  6 & 112     & 143     & $7.6 \times 10^8$   & \textbf{0} \\
  7 & 853     & 996     & $7.9 \times 10^9$   & \textbf{0} \\
  8 & 11,117  & 12,113  & $1.3 \times 10^{11}$ & \textbf{0} \\
  9 & 261,080 & 273,193 & $4.0 \times 10^{12}$ & \textbf{0} \\
  10 & 11,716,571 & 11,989,764 & 45\,s$^\dagger$ & \textbf{0} \\
  11 & 1,006,700,565 & 1,018,690,329 & 136\,min$^\dagger$ & \textbf{0} \\
  \bottomrule
\end{tabular}
\smallskip

\noindent{\footnotesize $^\dagger$For $m \le 9$, every tree on
$n \in \{5,\ldots,15\}$ vertices was tested (brute-force); the
Operations column counts hom evaluations. For $m = 10, 11$, only
near-path trees $E_n^{(d)}$ ($n \le 200$, $d \le 20$) were used
via the $E_n^{(d)}$ filter (wall-clock time shown). These rows are
bounded filter results in the sense of Remark~\ref{rem:filter}, not
full all-tree sweeps.}
\end{table}

\subsection{Asymptotic defense via Perron--Frobenius}

\begin{remark}[Filter heuristic]
  \label{rem:filter}
  The near-path family $E_n^{(d)}$ is a natural high-sensitivity filter for
  path-minimizer failures.
  The path $P_n$ achieves the minimal exponential growth rate
  $\lambda_1(H)^n$ (Csikv\'ari and Lin~\cite{csikvari2015}). For any
  connected, non-regular graph $H$, the Perron--Frobenius eigenvector equation guarantees
  $\Delta(H) > \lambda_1(H)$ strictly. Trees with macroscopic branching
  (such as stars) grow at an exponential rate strictly greater than $\lambda_1(H)^n$,
  exponentially faster than the path.
  To challenge the path, a tree must match the growth base $\lambda_1^n$
  while achieving a smaller leading coefficient; near-path trees
  $E_n^{(d)}$ provide a controlled way to test such localized boundary
  defects. We use this as a computational filter, not as an unconditional
  proof that all possible source-tree obstructions have been exhausted.
\end{remark}

This reduces the per-target cost from brute-force $O(t \cdot n \cdot m^2)$
(over $t$ source trees) to $O(d_{\max} \cdot n_{\max} \cdot m)$ via
$O(m)$ matrix--vector iteration up to $n = 200$, $d \le 20$.
Filtering across all 1,018,690,329 connected graphs with $m \le 11$
yields zero near-path violations in the tested range.
The exact exhaustive and certified-table computations in this paper use
integer arithmetic, so compiler flags such as \texttt{-march=native}
change runtime but not the certified values themselves; only
floating-point screens are treated as conditional filters before exact
verification.

\begin{table}[H]
\centering
\caption{Conditional results from the $E_n^{(d)}$ filter and bipartite sweeps.}
\label{tab:conditional}
\begin{tabular}{@{}lll@{}}
  \toprule
  Domain & Scope & Outcome \\
  \midrule
  All connected, $m \le 11$ & $n \le 200$, $d \le 20$ & no near-path hit \\
  Bipartite, $m \le 14$ & $n \le 199$, $d \le 20$ & no near-path hit \\
  Bipartite, $m = 15$ & $n \le 199$, $d \le 20$ & $H^*$ exists \\
  Bipartite depth-2, $m \le 17$ & $d=2$ sweep & no hit below $H_{18}$ \\
  Trees, $m \le 23$ & tested family & no hit \\
  Symmetric trees, $|V| \le 77$ & tested family & no hit \\
  \bottomrule
\end{tabular}
\end{table}

% ════════════════════════════════════════════════════════════════════════

\section{Computational Methodology}
\label{sec:computation}

Tree enumeration uses the Beyer--Hedetniemi canonical level-sequence
algorithm~\cite{BeyerH80}, generating all non-isomorphic trees on $n$ vertices in
$O(n)$ amortized time per tree. The \texttt{synthesizer} engine outputs
trees in graph6 format on \texttt{stdout}, which is piped to the
\texttt{leontovich\_fast} filter. The filter decodes graph6 strings,
precomputes $A^k \mathbf{1}$ for $k = 0, \ldots, n_{\max}$, and evaluates
$\Hom(E_n^{(d)}, H)$ vs.\ $\Hom(P_n, H)$ via matrix--vector iteration in
IEEE 754 double precision. Because the adjacency matrices and state vectors
are strictly non-negative, the iteration avoids catastrophic subtractive
cancellation. A standard forward-error estimate suggests an accumulated
relative rounding error on the order of
$n_{\max} \cdot m \cdot \varepsilon \approx 6.7 \times 10^{-13}$ (at
$m=15, n=200$, with machine epsilon $\varepsilon \approx 2.22 \times 10^{-16}$).
Comparisons use a relative tolerance of $10^{-11}$, which safely exceeds this
estimate by more than an order of magnitude. It is not an outward-rounded
error enclosure. This floating-point engine acts strictly as a high-speed
filter: all flagged anomalies are subsequently verified using exact
integer arithmetic. The filter can miss a true crossover whose relative margin
lies below the tolerance, so its negative results are conditional screens, not
unconditional nonexistence theorems.
The total cost per target graph is $O(n_{\max} \cdot m^2)$, giving
$O(h \cdot n_{\max} \cdot m^2)$ for the full sweep over $h$ targets.

Throughput: $\sim\!256{,}000$ graphs/second on six threads (Intel
Coffee Lake, AVX2 FMA) using dense $16 \times 16$ padded adjacency
matrices via optimized DAXPY loop
ordering and OpenMP parallelism. The bipartite sweep
of 608,930,559 connected bipartite graphs on $2 \le m \le 15$ vertices
(validated against OEIS A005142~\cite{oeis}) completes in approximately
35~minutes.

\begin{remark}[Depth-dependent bipartite Leontovich graph]
\label{rem:depth_anomaly}
While no $d = 2$ Leontovich violations were found, the $m = 15$
bipartite sweep revealed a single graph $H^*$ (Figure~\ref{fig:anomaly})
for which $\Hom(E_n^{(d)}, H^*) < \Hom(P_n, H^*)$ at several larger
depths, verified with exact integer arithmetic.
This bipartite graph on 15 vertices with partition sizes $(3,12)$
and degree sequence $[7, 5, 5, 2^5, 1^7]$ exhibits 300
such anomalies across $d \in \{14,16,18,20\}$ and odd
$n \le 199$; the first occurs at $(n,d)=(49,16)$ with
\[
  \Hom(P_{49},H^*)-\Hom(E_{49}^{(16)},H^*) =
  41{,}377{,}493{,}496{,}440{,}451{,}164 > 0.
\]
The maximum relative deficit in the tested range is
$\sim 7.6 \times 10^{-4}$.
No anomalies of any depth were \emph{flagged} at $m \le 14$ in this
floating-point screen; as above, this is conditional on the tolerance and does
not exclude an unflagged crossover with a smaller relative margin.
\end{remark}

\begin{figure}[H]
  \centering
  \includegraphics[width=0.38\textwidth]{anomaly_NCQCCA.png}%
  \hfill
  \includegraphics[width=0.58\textwidth]{anomaly_ratio.png}
  \caption{\textbf{Left:} The depth-anomaly graph $H^*$ on 15
    vertices\protect\footnotemark; red/blue shows the $(3,12)$ bipartition.
    \textbf{Right:} Ratio $\Hom(E_n^{(d)}, H^*) / \Hom(P_n, H^*)$ for
    odd~$n$. For the displayed larger depths the ratio drops below~1
    (shaded), meaning $E_n^{(d)}$ beats the path. For $d = 2$ the path
    remains the minimizer.}
  \label{fig:anomaly}
\end{figure}
\footnotetext{Graph6 encoding: \texttt{NCQCCA?\_B?K?W?g?K??}}

To search beyond the exhaustive enumeration wall,
we employ a simulated annealing search seeded from $T(7,1,9)$,
mutating edges and contracting vertices to seek smaller Leontovich
graphs within the graph manifold.

\begin{remark}
Since an independent set in $G$ is exactly a homomorphism
$G \to H_{\mathrm{ind}}$ (two vertices, exactly one of which has a loop, connected by an edge), the
framework developed here extends naturally to the Merrifield--Simmons
index: degree-constrained trees that maximize independent sets are
dual to trees that maximize homomorphisms into the same looped target.
\end{remark}

\section{SMT-Based Constraint Search and the 18-Vertex Depth-2 Discovery}
\label{sec:discussion}

The near-path sweeps for $m \le 11$ and the 15-vertex bipartite example
$H^*$ illustrate how quickly the Leontovich search space grows. Beyond the
range of exhaustive enumeration, we encode the homomorphism inequalities as
an integer constraint problem for a Satisfiability Modulo Theories (SMT)
solver. This search produces the 18-vertex depth-2 bipartite example
$H_{18}$ studied below.

\subsection{SMT-based algebraic search}

To bypass the exponential explosion of graph enumeration (e.g., more than $5 \times 10^{11}$ bipartite targets on $m = 16$ vertices),
the search for a Leontovich target graph $H$ of order $m$ can be formulated as a Constraint Satisfaction Problem (CSP)
over the binary adjacency variables $A_{ij} \in \{0, 1\}$.

Rather than explicitly unrolling the sum over all $m^n$ possible mappings to compute tree homomorphisms,
we embed the dynamic programming walk recurrences directly into the solver's constraint network as auxiliary variables.
For a target graph $H$, we define the walk vectors $w_k \in \mathbb{Z}^m$ recursively:
\begin{align}
  w_0(i) &= 1, \quad \forall i \in \{1, \dots, m\} \\
  w_k(i) &= \sum_{j=1}^{m} A_{ij} \, w_{k-1}(j), \quad \forall k \ge 1
\end{align}
The homomorphism count for the path is given by the linear sum
$\Hom(P_n, H) = \sum_{i=1}^{m} w_{n-1}(i)$. For the near-path family
$E_n^{(d)}$ with a pendant edge attached at depth $d$, the count can be
represented compactly as:
\begin{equation}
  \Hom(E_n^{(d)}, H) = \sum_{i=1}^{m} w_{n-d-2}(i) \, w_1(i) \, w_d(i)
\end{equation}
The search for a Leontovich target graph then reduces to solving the polynomial inequality:
\begin{equation}
  \label{eq:smt_ineq}
  \Hom(E_n^{(d)}, H) < \Hom(P_n, H)
\end{equation}
subject to adjacency constraints. Bipartiteness can be enforced dynamically
by introducing partition variables $x_i \in \{0, 1\}$ such that $A_{ij} = 1
\implies x_i \neq x_j$, while connectivity is guaranteed by defining
distance levels from a root vertex.

\subsection{Local SMT Leaf-Pruning Proof for \texorpdfstring{$T(7,1,9)$}{T(7,1,9)}}
\label{sec:local-pruning-audit}

To evaluate the optimality of our 76-vertex Leontovich graph
$T(7,1,[9^5, 8^2])$ discovered by simulated annealing, we formulated a
highly targeted local leaf-pruning search. Starting from the fully symmetric
tree $T(7,1,9)$ (78 vertices), we define integer decision variables $k_1,
\dots, k_7 \in \{0, \dots, 9\}$ representing the number of active leaves on
each of the 7 hubs. To break isomorphic permutations, we enforce $9 \ge k_1 \ge
k_2 \ge \dots \ge k_7 \ge 0$.

Because of the high symmetry of the quotient orbits, the walk equations can be
simplified from $m = 15 + \sum k_i$ variables down to a bilinear recurrence of
exactly $22$ variables per step:
\begin{align}
  w_k(\text{root}) &= \sum_{i=1}^{7} w_{k-1}(\text{bridge } i) \\
  w_k(\text{bridge } i) &= w_{k-1}(\text{root}) + w_{k-1}(\text{hub } i) \\
  w_k(\text{hub } i) &= w_{k-1}(\text{bridge } i) + k_i \, w_{k-1}(\text{leaf-type } i) \\
  w_k(\text{leaf-type } i) &= w_{k-1}(\text{hub } i)
\end{align}
We executed this local SMT solver to check if any pruning of size $|V| \le 75$ is
Leontovich for any threshold $n \le 31$. The solver returned \textbf{UNSAT} for every
single threshold.

This establishes a bounded computational audit:
\begin{compresult}
  Within the $T(7,1,9)$ star-of-stars pruning landscape, no pruning with
  $|V| \le 75$ is detected as Leontovich by the completed integer constraint
  checks for thresholds $n \le 31$. The 76-vertex tree
  $T(7,1,[9^5, 8^2])$ is therefore the minimum witness certified by this
  bounded audit.
\end{compresult}

\subsection{The 18-Vertex Depth-2 Bipartite Leontovich Graph Discovery}

To search for small non-tree Leontovich graphs, we restricted the
target class to bipartite graphs with left-partition size $m_1 \in \{2, 3\}$.
Under this partition constraint, every vertex in the right partition $R$ must
connect to a non-empty subset of the left partition $L$.

For $m_1 = 2$, there are only $2^2 - 1 = 3$ valid connection patterns for
vertices in $R$. A graph is thus uniquely specified up to isomorphism by a
triplet $(c_1, c_2, c_3)$ where $c_1 + c_2 + c_3 = m_2$.

\begin{proposition}[The $n=5$ depth-2 obstruction for $m_1 = 2$]
For every connected bipartite graph of partition $(2,m_2)$, the depth-2
near-path comparison has no crossover at $n=5$.
\end{proposition}
\begin{proof}
Let $H$ be a connected bipartite graph with left partition $L = \{u_1, u_2\}$.
The homomorphism counts at odd thresholds $n=2k+1$ are governed by the $2 \times 2$ transfer matrix $M = A^2|_L$, where $A$ is the adjacency matrix. We have:
\[
  M = \begin{pmatrix} c_1+c_3 & c_3 \\ c_3 & c_2+c_3 \end{pmatrix}.
\]
For $H$ to be connected, we must have $c_3 \ge 1$. Direct algebraic expansion of the transfer matrix product for $n=5$ yields the exact difference:
\[
  \Hom(P_5, H) - \Hom(E_5^{(2)}, H) = -c_3 \bigl( (c_1 - c_2)^2 + c_1 + c_2 \bigr).
\]
Since $c_1, c_2 \ge 0$ and $c_3 \ge 1$, this difference is non-positive, meaning $\Hom(P_5, H) \le \Hom(E_5^{(2)}, H)$.

\end{proof}

Beyond $n=5$, every $(2,m_2)$ target of total order $m\le15$ is covered at
every tested depth by the exhaustive bipartite sweep of
Section~\ref{sec:computation}, which found no crossover for odd $n\le199$ and
$d\le20$. We make no unconditional claim for larger $m_2$ or beyond that
window.

For $m_1 = 3$, there are $2^3 - 1 = 7$ valid connection patterns: 3
singletons ($c_1, c_2, c_3$), 3 pairs ($c_4, c_5, c_6$), and 1 triplet ($c_7$).
Any such bipartite target is uniquely determined by the vector $(c_1, \dots,
c_7)$.

The 15-vertex graph $H^*$ from Remark~\ref{rem:depth_anomaly} shows that
allowing arbitrary near-path depths yields a smaller bipartite Leontovich graph.
The following result isolates the smallest example found by the stricter
depth-2 test $E_n^{(2)}$, the near-path family used in the classical
Leontovich construction:
\begin{compresult}
  The bipartite graph $H_{18}$ of partition $(3, 15)$ with pattern vector
  $(7, 0, 0, 1, 1, 6, 0)$ is the unique minimal depth-2 bipartite Leontovich
  graph of partition $(3, m_2)$ found in this completed depth-2 bounded sweep
  through path length $n \le 51$. It first crosses against
  $E_n^{(2)}$ at $n = 17$. Its Graph6 representation is
  \begin{equation*}
    \texttt{QCaCCA?\_E?S?W?W?K?B??W?@\_??}
  \end{equation*}
\end{compresult}

Allowing arbitrary attachment depth gives an earlier crossover for the same
target. Exact integer arithmetic gives
\begin{align*}
  \Hom(P_{15},H_{18})&=1{,}105{,}256{,}666,\\
  \Hom(E_{15}^{(4)},H_{18})&=1{,}104{,}756{,}090,
\end{align*}
so the first near-path crossover for $H_{18}$ occurs at $(n,d)=(15,4)$,
with margin $500{,}576$. The symmetric depth $d=9$ gives the same source tree,
since $E_{15}^{(4)}\cong E_{15}^{(9)}$ by Remark~\ref{rem:d-symmetry}.
An exact scan over every admissible depth finds no crossover for $n<15$.

At the depth-2 crossover threshold $n=17$, the homomorphism counts verified by exact-precision arithmetic are:
\begin{align*}
  \Hom(P_{17}, H_{18}) &= 14{,}801{,}051{,}732 \\
  \Hom(E_{17}^{(2)}, H_{18}) &= 14{,}795{,}982{,}954
\end{align*}
yielding a positive Leontovich margin of $\Delta = 5{,}068{,}778$. Since
$m=18 < 76$, this is less than one quarter the order of the smallest
previously known depth-2 construction in the tested search regime. Because
the exhaustive check of all 45,838
non-isomorphic configurations for $m_2 < 15$ returned zero depth-2 hits,
\textbf{18 vertices is the depth-2 minimum within this partition family in the
completed depth-2 sweep through path length $n \le 51$}.

\subsubsection{\texorpdfstring{Structure and walk iteration of $H_{18}$}{Structure and walk iteration of H18}}

\begin{figure}[H]
  \centering
  \includegraphics[width=0.9\textwidth]{figures/leontovich_18.pdf}
  \caption{\texorpdfstring{The optimal 18-vertex depth-2 bipartite Leontovich graph $H_{18}$ discovered via SMT synthesis.}{The optimal 18-vertex depth-2 bipartite Leontovich graph H18 discovered via SMT synthesis.}}
  \label{fig:walk_h18}
\end{figure}

To make the mechanics of the Leontovich crossover transparent, we detail the
orbit decomposition and walk recurrences for $H_{18}$.
By partitioning the $18$ vertices into canonical symmetry groups, we reduce the
$18 \times 18$ transition adjacency representation into a $7 \times 7$ quotient system. The seven symmetry groups are:
\begin{enumerate}
  \item $L_0$: Left vertex $v_0$ ($1$ node, degree $9$).
  \item $L_1$: Left vertex $v_1$ ($1$ node, degree $7$).
  \item $L_2$: Left vertex $v_2$ ($1$ node, degree $7$).
  \item $R_{\mathrm{leaf},0}$: Right leaves connected only to $v_0$ ($7$ nodes, degree $1$).
  \item $R_{01}$: Cross-connector connected to $\{v_0, v_1\}$ ($1$ node, degree $2$).
  \item $R_{02}$: Cross-connector connected to $\{v_0, v_2\}$ ($1$ node, degree $2$).
  \item $R_{12}$: Bipartite core connected to $\{v_1, v_2\}$ ($6$ nodes, degree $2$).
\end{enumerate}

The step-by-step walk recurrence vectors $w_k(u)$ (denoting the number of walks of length $k$ starting at $u$)
are governed by the following exact coupled system:
\begin{align*}
  w_k(L_0) &= 7 w_{k-1}(R_{\mathrm{leaf},0}) + w_{k-1}(R_{01}) + w_{k-1}(R_{02}) \\
  w_k(L_1) &= w_{k-1}(R_{01}) + 6 w_{k-1}(R_{12}) \\
  w_k(L_2) &= w_{k-1}(R_{02}) + 6 w_{k-1}(R_{12}) \\
  w_k(R_{\mathrm{leaf},0}) &= w_{k-1}(L_0) \\
  w_k(R_{01}) &= w_{k-1}(L_0) + w_{k-1}(L_1) \\
  w_k(R_{02}) &= w_{k-1}(L_0) + w_{k-1}(L_2) \\
  w_k(R_{12}) &= w_{k-1}(L_1) + w_{k-1}(L_2)
\end{align*}
with the base canonical condition $w_0(u) = 1$ for all vertices. The explicit numerical values for the
first five steps are shown in Table~\ref{tab:walk_h18}. This iteration can be carried out exactly to $w_{16}(u)$.

\begin{table}[H]
\centering
\caption{Step-by-step walk vector iteration and total sum counts for $H_{18}$.}
\label{tab:walk_h18}
\begin{tabular}{lcccccccc}
  \toprule
  Step $k$ & $L_0$ & $L_1$ & $L_2$ & $R_{\mathrm{leaf},0}$ ($\times 7$) & $R_{01}$ ($\times 1$) & $R_{02}$ ($\times 1$) & $R_{12}$ ($\times 6$) & Total $\sum w_k$ \\
  \midrule
  $w_0$ & 1 & 1 & 1 & 1 & 1 & 1 & 1 & 18 \\
  $w_1$ & 9 & 7 & 7 & 1 & 2 & 2 & 2 & 46 \\
  $w_2$ & 11 & 14 & 14 & 9 & 16 & 16 & 14 & 218 \\
  $w_3$ & 95 & 100 & 100 & 11 & 25 & 25 & 28 & 590 \\
  $w_4$ & 127 & 193 & 193 & 95 & 195 & 195 & 200 & 2{,}768 \\
  $w_5$ & 1055 & 1395 & 1395 & 127 & 320 & 320 & 386 & 7{,}690 \\
  \bottomrule
\end{tabular}
\end{table}

The path count is given by the sum $\Hom(P_{17}, H_{18}) = \sum_u w_{16}(u)$, while the near-path count
is computed as $\Hom(E_{17}^{(2)}, H_{18}) = \sum_u w_{13}(u) w_1(u) w_2(u)$, yielding the exact margin reported above.

\begin{remark}[Rigorous Resolution of Other Bipartite Families]
  To determine whether smaller depth-2 examples occur in other bipartite
  partition families, we engineered a high-performance C++ solver to
  exhaustively sweep all other partition families $(m_1, m_2)$ for
  $m_1 \ge 4$ up to total order $m_1 + m_2 \le 17$. By evaluating all
  non-isomorphic bipartite connection patterns up to path length $n = 51$
  (covering millions of distinct configurations), the solver returned zero
  depth-2 hits. This supports the claim that $H_{18}$ is the smallest
  currently verified depth-2 bipartite Leontovich graph across the tested
  partition families in the completed depth-2 sweep through path length
  $n \le 51$; beyond that bounded regime, Section~\ref{sec:doublecover}
  shows why larger thresholds can behave differently. The depth-dependent
  graph $H^*$ remains the smaller bipartite Leontovich example for
  general $E_n^{(d)}$.
\end{remark}

\subsection{The Near-Miss Asymptotic Dilution}

During our trace of the $m_1 = 3$ family, we discovered pattern vectors where
the homomorphism count ratio approached $1.0$ from above with extreme
precision (e.g., the 14-vertex target $c = (9, 1, 0, 0, 0, 1, 0)$).
However, algebraic inspection reveals these are not impending crossovers, but
disjoint unions.

The vector $(9, 1, 0, 0, 0, 1, 0)$ specifies a disconnected graph
$H \cong K_{1,9} \sqcup P_4$. Because homomorphism counts map connected trees
entirely into a single connected component, they are strictly additive over
disjoint components: $\Hom(T, H) = \Hom(T, K_{1,9}) + \Hom(T, P_4)$. The star
$K_{1,9}$ is degenerate (Proposition~\ref{prop:p3}) and yields
identical counts for $P_n$ and $E_n^{(2)}$, contributing a massive
$O(3^n)$ perfect tie term. The path $P_4$ provides a strict positive margin,
but at a much slower growth rate $O(1.618^n)$. The ratio approaches $1.0$
precisely because it is physically impossible to cross over; the strict minimality of $P_4$
is asymptotically diluted by the dominant degenerate star.

\subsection{The parameterized \texorpdfstring{$\mathcal{B}(x,y)$}{B(x,y)} family}

The pattern $(7,0,0,1,1,6,0)$ of $H_{18}$ belongs to a natural
two-parameter family of bipartite targets.

Let $\mathcal{B}(x,y)$ be the bipartite graph defined by the left partition
$L = \{u_0, u_1, u_2\}$ and the connection pattern $(x, 0, 0, 1, 1, y, 0)$.
This graph has left partition size $m_1 = 3$ and right partition size
$m_2 = x + y + 2$, giving a total order of $m = x + y + 5$.

Due to the symmetry of the connection subsets, $u_1$ and $u_2$ form a perfect
automorphism orbit. The random walks restricted to the left partition taking
steps of length 2 are governed by a $3 \times 3$ matrix. However, the
antisymmetric eigenvector $(0, 1, -1)^T$ has an exact eigenvalue of $\lambda^2=1$,
but is perfectly orthogonal to the all-ones initial state vector $\mathbf{1}$.
Thus, it contributes nothing to the homomorphism counts. The active dynamics
collapse into the $2 \times 2$ asymmetric quotient matrix:
\begin{equation}
  M = A^2|_L = \begin{pmatrix} x+2 & 2 \\ 1 & 2y+1 \end{pmatrix}.
\end{equation}

The non-zero squared eigenvalues of $\mathcal{B}(x,y)$ are precisely the roots
of the characteristic polynomial of $M$:
\begin{equation}
  \lambda^2 - (x + 2y + 3)\lambda + (2xy + x + 4y) = 0.
\end{equation}

For our record-breaking graph $H_{18}$, we evaluate $\mathcal{B}(7,6)$. The
characteristic polynomial simplifies to $\lambda^2 - 22\lambda + 115 = 0$,
yielding exact squared eigenvalues of $11 \pm \sqrt{6}$. The active eigenvalue
ratio governing the crossover is thus exactly:
\begin{equation}
  r = \sqrt{\frac{11 - \sqrt{6}}{11 + \sqrt{6}}} \approx 0.797.
\end{equation}

The discriminant $\Delta = (x-2y+1)^2 + 8$ shows that the
ratio is maximized precisely along the ``resonance line'' $x = 2y - 1$,
which favors long-lived eigenvalue interference. Computations suggest that
this parameterized family contains infinitely many Leontovich graphs, but we
do not prove that assertion here. Within the completed
depth-2 $(3,m_2)$ block-pattern sweep through path length $n \le 51$,
$H_{18}$ is the smallest member of $\mathcal{B}(x,y)$ found to cross over.

% ════════════════════════════════════════════════════════════════════════

\section{The Bipartite Double Cover and Strongly Leontovich Graphs}
\label{sec:doublecover}

A natural way to remove a loop from a target is to pass to the
bipartite double cover; Galvin et al.~\cite{galvin2026} note the
equivalent two-copy construction. We use this language to state the
preservation result and the even crossover threshold for the resulting
simple bipartite graph.

\begin{definition}[Bipartite Double Cover]
For any graph $H$, its bipartite double cover $B = H \times K_2$ is a
bipartite graph on $2|V(H)|$ vertices. The vertex set is $V(H) \times \{0, 1\}$,
with edges connecting $(u, 0)$ to $(v, 1)$ if and only if $(u, v) \in E(H)$.
A loop at vertex $u$ in $H$ corresponds to a standard edge between $(u, 0)$ and
$(u, 1)$ in $B$. Thus, $B$ is always a simple bipartite graph.
\end{definition}

\begin{theorem}
\label{thm:double_cover}
A graph $H$ is Leontovich if and only if its bipartite double cover $B = H
\times K_2$ is Leontovich. Furthermore, $H$ is Strongly Leontovich if and only
if $B$ is Strongly Leontovich.
\end{theorem}
\begin{proof}
Let $T$ be any connected tree. Because $T$ is bipartite, it admits exactly two
proper 2-colorings. Consequently, any homomorphism $\phi \colon T \to H$ lifts
to exactly two distinct homomorphisms $\Phi \colon T \to B$, corresponding to
the two choices of mapping the color classes of $T$ to the partitions $\{0, 1\}$
of $B$. Therefore, for any tree $T$, $\Hom(T, B) = 2 \Hom(T, H)$.

It follows immediately that for all $n$ and $d$:
\[
  \Hom(P_n, B) - \Hom(E_n^{(d)}, B) = 2 \Bigl[ \Hom(P_n, H) - \Hom(E_n^{(d)}, H) \Bigr]
\]
Because the homomorphism difference scales by a strict positive constant, the
crossover thresholds and the asymptotic ratio of the leading projection
coefficients $c_1(E_n)/c_1(P_n)$ are perfectly preserved.
\end{proof}

\subsection{A Simple Strongly Leontovich Graph}

The coefficient calculation in this section establishes that the looped symmetric
tree $\hat{T}(1, 35, 1, 50)$ on 1,822 vertices is Strongly Leontovich, with an asymptotic leading
coefficient ratio of

$\lim_{n \to \infty} \Hom(E_n^{(2)}, \hat{T}) / \Hom(P_n, \hat{T})
\approx 0.999953714414 < 1$.
Exact finite counts at $n=15001$ give the
still-unconverged ratio $0.999864892078$, reflecting the very slow decay of the
near-antipodal mode at eigenvalue $-7.282828$, whose magnitude is within a
factor $0.9997162$ of the Perron root $7.284896$.

By Theorem~\ref{thm:double_cover}, its bipartite double cover $B = \hat{T}
\times K_2$ inherits this exact property. $B$ is a simple, loopless, bipartite
graph on exactly 3,644 vertices. Because its asymptotic ratio is strictly less
than 1, $B$ is a \textbf{simple strongly Leontovich graph}, addressing the
simple-target existence question.

Furthermore, because $B$ is Strongly Leontovich, $\Hom(E_n, B) < \Hom(P_n, B)$
for all sufficiently large $n$, both even and odd. Exact integer computation of
the walk sequences reveals that $E_n$ overtakes $P_n$ on $B$ starting exactly at
the even threshold $n = 17{,}340$ and remains ahead through the verified range.
So even-$n$ crossovers do not require loops. Since $B$ is bipartite its
spectrum is exactly symmetric about zero, and that $\pm$ pairing is what
protects $P_n$ at even~$n$. What eventually defeats the protection is
near-degeneracy among the dominant magnitudes: the two largest positive
eigenvalues of $B$ are $7.284896$ and $7.282828$, a ratio of about
$0.99972$, so the competing mode decays slowly enough to survive until
$n = 17{,}340$. Equivalently, on the looped source $\hat{T}(1,35,1,50)$
the same fact appears as a near-antipodal eigenvalue $-7.282828$ just
inside the Perron root. This is not the second positive eigenvalue of
$\hat{T}$, $\lambda_2 \approx 5.842228$: at
$\lambda_2/\lambda_1 \approx 0.802$ that mode is negligible by
$n \approx 45$.

\paragraph{Disproof of the universal bipartite parity hypothesis.}
Appendix~\ref{app:sa} records a working hypothesis suggested by
extensive finite search: the 608,930,559 connected bipartite graphs on
$m \le 15$ vertices, together with all source trees through $n=22$
(including $5{,}623{,}756$ at $n=22$), showed no even-$n$ crossover. The
$n=17{,}340$ example above disproves the universal statement. The point
is scale, not a new mechanism: the relevant long-lived term is the
second dominant magnitude $7.282828$, not the smaller positive
eigenvalue $\lambda_2$.

\subsection{A Perturbed Non-Bipartite Certificate}

Adding a single edge inside one side of a bipartite strongly Leontovich
graph can produce a simple non-bipartite target. Here we use the looped
symmetric tree $\hat{T}(2, 34, 1, 48)$ on 3,403 vertices, whose
bipartite double cover $B' = \hat{T}(2, 34, 1, 48) \times K_2$ is a
simple bipartite strongly Leontovich graph on 6,806 vertices with
leading-coefficient ratio $0.999713000053\ldots$ and margin
$2.87 \times 10^{-4}$.

Adding a single edge inside one side of this bipartite graph creates an odd
cycle, so the perturbed target is strictly non-bipartite. We use
$\hat{T}(2, 34, 1, 48)$ rather than the smaller $\hat{T}(1,35,1,50)$ because
the larger asymptotic margin survives the perturbation more comfortably.
Adding an edge $e$ between two sibling terminal leaves of $B'$ gives the
6,806-vertex graph $B' + e$. The script
\texttt{scripts/verify\_perturbed\_cert.py} computes the coarsest equitable
partition by colour refinement, verifies equitability vertex-by-vertex with
exact integer arithmetic, and obtains an 18-cell quotient matrix. On this
quotient, the leading-coefficient ratio is
\[
  \lim_{n \to \infty}
  \frac{\Hom(E_n^{(2)}, B' + e)}{\Hom(P_n, B' + e)}
  \approx 0.999952140676 < 1.
\]
The same pipeline reproduces the unperturbed ratio
$0.999713000053\ldots$ for $B'$. A Rayleigh quotient estimate explains
why the perturbation is stable: the principal eigenvector mass on the
added leaves is tiny, and the measured shift is
$\Delta\lambda_1 \approx 1.77 \times 10^{-4}$.

To estimate the bound, \texttt{scripts/strong\_coeff.py}
(\texttt{certified\_leading\_ratio}) computes a Bauer--Fike residual
and a Davis--Kahan angle bound for the symmetrized 18-cell quotient,
both at 50 digits. Denote the resulting numerical estimate by
$\rho_{\mathrm{hi}}$. Numerically,
$\|r\|_2 \approx 7.5\times 10^{-50}$ against a spectral gap
$\lambda_1 - \lambda_2 \approx 5.1\times 10^{-3}$, propagating to
$\rho_{\mathrm{hi}} - \rho_2(B'+e) \approx 7.5\times 10^{-47}$ and
\[
  \rho_{\mathrm{hi}} \approx 0.99995214067566 < 1,
  \qquad 1 - \rho_{\mathrm{hi}} \approx 4.79\times 10^{-5}.
\]
The estimated error is roughly 41 orders of magnitude smaller than the
margin, so the numerical estimate $\rho_{\mathrm{hi}}$ stays below $1$. This is
high-precision numerical evidence (N) for a simple non-bipartite
strongly-Leontovich candidate in the leading-coefficient sense. It is a
residual-based estimate in that sense, not a directed-rounding
interval-arithmetic certificate.

% ════════════════════════════════════════════════════════════════════════

\section{Conclusion and Open Problems}
\label{sec:conclusion}

Quotient-matrix reductions and exact computational sweeps constrain the
landscape of Leontovich graphs sharply. Proposition~\ref{prop:n5} rules
out $n=5$ in the spherically symmetric family $T(x,y,z)$. The smallest
crossover threshold witnessed in the
families studied here is $n=7$; two universal $n=6$ polynomial
inequalities remain open.

We conclude with the two principal open directions highlighted by this work:

\begin{enumerate}
  \item \textbf{Near-path sufficiency.} Every known Leontovich graph is
    detected by the near-path family $E_n^{(d)}$. The near-path
    sufficiency conjecture asks whether this always holds. If it does,
    the bounded near-path sweeps in this paper become bounded all-tree
    results over the same tested target orders and source-size windows;
    removing those finite bounds would still require additional arguments.
  \item \textbf{Minimal simple targets.} The bipartite double cover of
    $\hat{T}(1,35,1,50)$ yields a simple strongly Leontovich graph on
    3,644 vertices, and the perturbed quotient
    (Section~\ref{sec:doublecover}) gives high-precision numerical evidence
    for a simple non-bipartite candidate on 6,806 vertices. Closing the
    gap to the 11-vertex frontier probed by the simple-graph sweep
    remains a computational challenge.
\end{enumerate}

\section*{Code and Data Availability}
The C++ search routines, exact integer arithmetic formulations, and SMT
solver scripts supporting the exact finite computations \textbf{(E)}
and bounded sweeps \textbf{(S)} are available in the project
repository. The
\href{https://github.com/sjjaroch/leontovich-utils/tree/7aab6dd68bccd7614794c75333dddc692b53541f}{archival fork at commit \texttt{7aab6dd68bcc}}
is a mirror of the project repository pinned to the exact shared-history
snapshot used for the reported computations. An experimental Lean~4 scaffold
records intended theorem statements; no principal result is formally
verified, and the scaffold should not be read as verification. The executable
witness verifiers and experimental scaffold are maintained at
\url{https://github.com/sjjaroch/leontovich-utils}.

\section*{Acknowledgments}

The author thanks JD Nir for helpful correspondence and David Galvin
for the bipartite double-cover idea, both of which furthered this work.

\appendix
\input{sections/11_threshold_landscape}
\section{Simulated Annealing Results}
\label{app:sa}

Starting from $T(7,1,9)$ ($|V|=78$), our simulated annealing
search discovered a smaller Leontovich graph with $|V|=76$ via
edge mutation and vertex contraction.  The result was verified
independently with exact integer arithmetic
(\texttt{scripts/verify\_76.py}).

The 76-vertex graph is a tree with 75 edges and degree sequence
\[
  [10, 10, 10, 10, 10, 9, 9, 7, \underbrace{2, \ldots, 2}_{7},
   \underbrace{1, \ldots, 1}_{61}].
\]
It retains the star-of-stars topology of $T(7,1,9)$ but with
asymmetric leaf counts at the bridge layer (two branches each lost one
leaf).

\begin{table}[H]
\centering
\caption{Active quotient spectra along the pruning sequence. Asymmetric pruning breaks identical-branch symmetry, expanding the quotient matrix and lifting new active eigenvalues $\lambda_2^*$ that govern the Leontovich crossover.}
\label{tab:sa_hits}
\begin{tabular}{@{}clccccc@{}}
  \toprule
  $|V|$ & Structure & Crossover $n$ & $d$ & $\lambda_2^*/\lambda_1$ & $\lambda_1$ & $\lambda_2^*$ \\
  \midrule
  78 & $T(7,1,[9^7])$     & 13  & 2 & 0.6877 & 3.3973 & 2.3364 \\
  77 & $T(7,1,[9^6,8^1])$ & 15  & 2 & 0.8911 & 3.3874 & 3.0187 \\
  76 & $T(7,1,[9^5,8^2])$ & 17  & 2 & 0.8997 & 3.3766 & 3.0380 \\
  \midrule
  75 & $T(7,1,[9^4,8^3])$ & --- & --- & 0.9089 & 3.3647 & 3.0581 \\
  \bottomrule
\end{tabular}
\end{table}

\noindent
The 76-vertex graph has an active secondary quotient eigenvalue of $\lambda_2^* \approx 3.0380$. Its first crossover
at $n = 17$ with $d = 2$ was verified by exact integer computation
of $\Hom(P_n, H) - \Hom(E_n^{(2)}, H) = 1{,}993{,}646 > 0$.
For $n \ge 17$, the sign of $\Delta_n$ strictly alternates between
even and odd $n$, indicating destructive interference from the
active secondary eigenvalue.

A 20-seed swarm of $5 \times 10^6$ SA steps each (totalling
$10^8$ evaluations) was run from $T(7,1,9)$ at temperature 3.0.
Every seed converged to a 76-vertex graph, with an intermediate
$|V|=77$ Leontovich graph consistently discovered along the path
(typically within the first 10 steps). No seed found a graph with
$|V| < 76$.  This strongly suggests that 76 is the global minimum
within the star-of-stars basin, and that the descent
$78 \to 77 \to 76$ represents the unique pruning sequence.

\begin{figure}[H]
\centering
\includegraphics[width=0.75\textwidth]{figures/leontovich_76.pdf}
  \caption{The 76-vertex Leontovich graph $T(7,1,[9^5,8^2])$. Red: root
  (degree~7). Blue: bridges (degree~2). Orange: hubs (degree~9--10).
  Green: leaves (degree~1). The asymmetric pruning of one leaf from
  each of two branches in $T(7,1,9)$ expands the active quotient
  spectrum, shifting the first crossover from $n=13$ to $n=17$.}
\label{fig:leo76}
\end{figure}

\subsection*{Symmetry breaking and the active quotient spectrum}

The pruning sequence admits a clean analytic
description that resolves the crossover behaviour. In the fully symmetric
graph $T(7,1,[9^7])$, the adjacency matrix possesses highly degenerate
eigenvalues at $\lambda = \sqrt{10}$ (arising from the isolated 9-leaf hubs).
However, the eigenvectors associated with this eigenspace are antisymmetric
across the identical branches. Consequently, they are perfectly orthogonal
to the all-ones vector $\mathbf{1}$ and the initial degree vector, meaning
their projection coefficients in the transfer matrix expansion
$\Hom(G_n, H) = \mathbf{1}^T A^{n-4} \mathbf{h}_0$ are identically zero.
These orthogonal eigenvectors contribute nothing to the homomorphism counts.

Instead, the crossover dynamics are governed exclusively by the \emph{active}
spectrum---the eigenvalues of the quotient matrix over the graph's
automorphism orbits. For the fully symmetric graph $T(7,1,[9^7])$, the
active secondary eigenvalue is $\lambda_2^* \approx 2.3364$.

When leaves are pruned asymmetrically, the perfect identical-branch symmetry
is broken. For $T(7,1,[9^k, 8^{7-k}])$, the orbits split and the quotient
matrix expands from $4 \times 4$ to $7 \times 7$. This symmetry breaking
activates a new secondary eigenvalue $\lambda_2^*$ that directly interferes
with the Perron mode.

The Leontovich crossover arises from destructive interference between the
dominant Perron mode $\lambda_1^n$ and the secondary active mode $(\lambda_2^*)^n$
in the difference $\Hom(P_n, H) - \Hom(E_n^{(d)}, H)$. As branches are pruned,
$\lambda_1$ decreases while the active secondary eigenvalue $\lambda_2^*$
increases (see Table~\ref{tab:sa_hits}), narrowing the spectral gap.
At $|V| = 76$, the interference from $\lambda_2^* \approx 3.0380$ is just
sufficient to drive the ratio $\Hom(E_n^{(2)}, H)/\Hom(P_n, H)$ below~1 at
$n = 17$. The local SMT experiment for $|V| = 75$ times out at its first
depth-2 threshold, $n=13$, so it does not establish minimality within this
pruning family.

This yields a clean dichotomy: \emph{within the star-of-stars basin,
asymmetric pruning breaks orbital symmetry and lifts new active eigenvalues
into the quotient spectrum.} The Leontovich property is lost when this
symmetry-breaking pushes the active secondary eigenvalue too close to the
spectral radius. The 76-vertex graph is the smallest witness found by the
search, not a certified minimum; Section~\ref{sec:local-pruning-audit}
describes the incomplete local audit.

\subsection*{Deeper symmetric trees}

An overflow-safe sweep (\texttt{\_\_int128} exact arithmetic with
\texttt{\_\_builtin\_mul\_overflow}) over 596,891 branching sequences
$T(d_1, \ldots, d_k)$ with branching depth $k \le 5$ (at most six
orbits) and $|V| \le 2{,}000$
found 7,016 Leontovich instances and confirms that $T(7,1,9)$ at
$|V| = 78$ is the vertex-count minimum within the swept symmetric
branching-sequence family (depth $k\le5$, $|V|\le2{,}000$). For each
threshold~$n$, the smallest 5-orbit tree is
consistently \emph{larger} than the corresponding 4-orbit tree:
for example, at $n = 13$, $T(1,8,1,12)$ ($|V| = 114$) versus
$T(7,1,9)$ ($|V| = 78$).  Deeper trees do extend the reachable
threshold range---5-orbit entries reach $n \ge 61$ where no 4-orbit
tree appears---but at vertex counts exceeding 500.

An earlier floating-point sweep produced false positives (e.g.,
$T(1,1,1,21)$ at $|V| = 25$ was incorrectly flagged as Leontovich).
Exact integer verification showed $\Hom(E_n, H) > \Hom(P_n, H)$ at
all~$n$ for these candidates.  The root cause was signed
\texttt{\_\_int128} overflow (undefined behaviour under \texttt{-O3}),
which silently corrupted the comparisons; the
\texttt{\_\_builtin\_mul\_overflow} intrinsic eliminates this class
of error.

\subsection*{The Diophantine threshold frontier}

The minimum vertex counts in Table~\ref{tab:landscape_full} are
explained algebraically and certified computationally in the tested range.

The audited frontier is restricted to the subfamily $T(x,1,z)$. It does
not establish that contracting an arbitrary $T(x,y,z)$ with $y>1$
preserves or advances its crossover threshold. Within the audited
subfamily, the exact landscape sweep checked all $134{,}391{,}973$
graphs with $|V| \le 10{,}000{,}000$ against $E_n^{(2)}$ for odd
$n \le 51$ using exact integer arithmetic.

\begin{compresult}[Frontier certification in the tested range]
\label{thm:frontier}
For $T(x,1,z)$, the difference $\Delta_n(x,z) = \Hom(P_n, T) -
\Hom(E_n^{(2)}, T)$ is an exact polynomial in~$x$ and~$z$.
Within the tested range of Table~\ref{tab:landscape_full}, the
minimum-vertex graph achieving first crossover at odd threshold~$n$
is certified by exhaustive exact search over the integer pairs
$(x,z)$ satisfying $\Delta_n(x,z) > 0$ and $\Delta_{n-2}(x,z) \le 0$.
The dominant cubic of $\Delta_n$ explains the geometric shape of the
feasible region but is not, by itself, a complete proof of minimality
at the small thresholds realised here.
\end{compresult}

\begin{proof}[Derivation of the feasible wedge]
Both $\Hom(P_n, T)$ and $\Hom(E_n^{(2)}, T)$ are computed via
finite matrix multiplications over the $4 \times 4$ quotient matrix.
Their difference is therefore an exact polynomial.  For example:
\begin{align*}
  \Delta_7(x,z) &= -x\bigl(x^3 - x^2 - 2xz^2 - xz - x + (z+1)^3\bigr).
\end{align*}
Setting $\Delta_7 > 0$ and extracting the dominant cubic,
$-x^3 + 2xz^2 - z^3 > 0$, the substitution $w = x/z$ yields the
characteristic equation $-w^3 + 2w - 1 = 0$, which factors as
$-(w-1)(w^2 + w - 1) = 0$.  The positive roots give
\[
  \frac{\sqrt{5} - 1}{2} < \frac{x}{z} < 1,
\]
i.e., the dominant terms place the feasible integer pairs inside a
geometric wedge bounded by the inverse golden ratio
$\varphi^{-1} \approx 0.618$. This explains the observed frontier
shape and sharply narrows the exact search region, but the final
minimality claims below are certified by exhaustive integer search in
the tested range rather than by the wedge alone. For example, at
$n = 7$ the
unique minimum is $(x,z) = (16,19)$ with $\Delta_7(16,19) = 512 > 0$
and $\Delta_5(16,19) = -560 \le 0$, yielding $|V| = 337$.  This
improves the $1{,}801$-vertex bound of
Galvin et al.~\cite{galvin2025} to 337 vertices. Exact exhaustive search then certifies the remaining entries in
Table~\ref{tab:landscape_full}: $(9,11)$ at $n = 9$; $(8,10)$ at
$n = 11$; $(7,9)$ at $n = 13$; and $(8,12)$ at $n = 15$, matching
the computational landscape in Table~\ref{tab:landscape_full}.
\end{proof}

\begin{remark}[Algebraic ballast in the strongly Leontovich case]
\label{rem:ballast}
A graph is \emph{strongly Leontovich}~\cite{galvin2026} if
$E_n$ beats the path asymptotically, i.e.\ the leading coefficient ratio
satisfies $\rho_d(H) < 1$ in Eq.~\eqref{eq:rho}.  This is harder than
the standard Leontovich condition, which only requires destructive
interference at \emph{some} odd~$n$.

Under 4-orbit symmetry, the quotient dynamics are governed by a
$4 \times 4$ quotient matrix with only 3 integer degrees of freedom
$(x, y, z)$.  To force the strict ratio inequality $\rho_d(H) < 1$ in
Eq.~\eqref{eq:rho} within this tiny parameter space, the branching degrees
must be scaled to extreme values (e.g., $x = 400$, $z = 800$ for
$\hat{T}(400, 3, 800)$).  The resulting $\sim\!960{,}000$ vertices are
\emph{algebraic ballast}: redundant, identical branches imposed by the
symmetry constraint, carrying eigenvectors orthogonal to the initial state
that project to zero in the homomorphism counts and contribute nothing to
the crossover.

Adding a 5th orbit expands the quotient matrix from $4 \times 4$ to
$5 \times 5$, unlocking a new independent variable.  This produces many
small ordinary Leontovich examples, but the strong condition must be checked
asymptotically rather than by observing a finite crossover window.  A corrected
exact audit of 5-orbit looped symmetric trees
$\hat{T}(d_1, d_2, d_3, d_4)$ finds that the smallest verified
\textit{strongly} Leontovich example in this sweep is
$\hat{T}(1, 35, 1, 50)$ on $1{,}822$ vertices.  This still improves the
$961{,}601$-vertex 4-orbit upper bound by roughly a factor of $528$, while the
smaller examples in Table~\ref{tab:strong_frontier} are ordinary finite-window
Leontovich graphs rather than strongly Leontovich graphs.
For every row in the table, the sign of $\rho_2-1$ is additionally certified
over the exact algebraic Perron root: the verifier isolates that root in a
rational interval and proves that the corresponding integer sign polynomial
has no zero in the interval. Thus these asymptotic classifications do not
depend on a floating-point tolerance.

\begin{table}[H]
\centering
\caption{Audit of 5-orbit looped symmetric trees: ordinary versus strongly Leontovich behavior.}
\label{tab:strong_frontier}
\begin{tabular}{@{}ccll@{}}
\toprule
Opening $n$ & $|V|$ & Parameters $(d_1, d_2, d_3, d_4)$ & Verdict \\
\midrule
7 & 1,066 & $\hat{T}(1, 28, 1, 36)$ & ordinary only; closes at $n=1827$ \\
9 & 370 & $\hat{T}(1, 16, 1, 21)$ & ordinary only; closes at $n=251$ \\
11 & 262 & $\hat{T}(1, 13, 1, 18)$ & ordinary only; closes at $n=141$ \\
13 & 242 & $\hat{T}(1, 12, 1, 18)$ & ordinary only; closes at $n=127$ \\
15 & 546 & $\hat{T}(1, 17, 1, 30)$ & ordinary only; closes at $n=261$ \\
17 & 422 & $\hat{T}(1, 15, 1, 26)$ & ordinary only; closes at $n=133$ \\
\midrule
7 & 1,822 & $\hat{T}(1, 35, 1, 50)$ & strongly Leontovich (smallest known record) \\
\bottomrule
\end{tabular}
\end{table}
\end{remark}

\subsection*{Non-oscillatory threshold}

\begin{proposition}
Suppose the odd-$n$ difference has the exact two-mode form
\[
  \Delta_n = \Hom(P_n, H) - \Hom(E_n^{(d)}, H)
  = A\,\lambda_1^n + B\,\lambda_2^n
\]
for some $A>0$, $B<0$, and $0 < \lambda_2 < \lambda_1$. Then
$\Delta_n$ restricted to odd~$n$ can change sign at most once. This
statement does not apply to the 5-orbit looped symmetric targets in
Table~\ref{tab:strong_frontier}.
\end{proposition}

\begin{proof}[Proof sketch]
Write $r = \lambda_2/\lambda_1 \in (0,1)$. Then
\[
  \Delta_n = \lambda_1^n\bigl(A + B r^n\bigr)
\]
and $B r^n$ is strictly increasing because $B<0$ and $r^n$ decreases
with $n$. Hence $A + B r^n$ crosses zero at most once on the odd
subsequence. Any finite base cases still need to be checked directly.
\end{proof}

\noindent
We verified this computationally for the 76-vertex graph: $\Delta$
is negative at odd $n \le 15$ and strictly positive at odd
$17\le n\le43$, with the gap growing
exponentially from $\Delta_{17} = 1{,}993{,}646$ to
$\Delta_{43} > 10^{17}$.

\subsection*{A Bipartite Parity Hypothesis}

\begin{remark}[Bipartite parity hypothesis (refuted in Section~\ref{sec:doublecover})]
For any bipartite target graph $H$, $\Hom(E_n^{(d)}, H) \ge
\Hom(P_n, H)$ for all even~$n$.  That is, the path $P_n$ is
always the minimizer among trees on $n$ vertices when $n$ is even
and $H$ is bipartite.
\end{remark}

\noindent
Extensive finite searches suggested this hypothesis: across all tested
target graphs (the 608,930,559 connected bipartite graphs on $m \le 15$
vertices, together with all source trees through $n=22$, including
$5{,}623{,}756$ at $n=22$), no even-$n$ crossover was observed.
Section~\ref{sec:doublecover}
disproves it: the bipartite double cover $B$ of $\hat T(1,35,1,50)$
overtakes the path at the even threshold $n = 17{,}340$.

\begin{remark}
The same parity effect appears for independent sets: the consecutive
ratios $a(N)/a(N-1)$ oscillate with period~2; the looped target
$H_{\mathrm{ind}}$ breaks bipartiteness, so the parity protection is
gone.
\end{remark}

\bibliographystyle{plain}
\bibliography{references}

\end{document}
